\documentclass[12pt,a4paper]{article}
\usepackage[utf8]{inputenc}
\usepackage[T1]{fontenc}
\usepackage{amsmath,amssymb,amsfonts}
\usepackage{hyperref}
\usepackage{geometry}
\usepackage{booktabs}
\usepackage{amsthm}

\newtheorem{theorem}{Theorem}[section]
\newtheorem{lemma}[theorem]{Lemma}
\newtheorem{proposition}[theorem]{Proposition}
\newtheorem{corollary}[theorem]{Corollary}
\newtheorem{definition}[theorem]{Definition}
\newtheorem{remark}[theorem]{Remark}
\newtheorem{assumption}[theorem]{Assumption}
\newtheorem{observation}[theorem]{Observation}

\geometry{a4paper, margin=1in}

\title{Crystals Cannot Be Grown: A Large-Deviation Theory of Typical and Rigid Collatz Structures}
\author{Sergey Alexandrovich Zadorozhniy}
\date{August 2026}

\begin{document}
\maketitle

\begin{abstract}
We develop a unified large-deviation theory of the Collatz (Syracuse) space that
explains the dichotomy between \emph{typical decay} and \emph{rigid confluence
structures}, and we show that the two are two faces of a single rate function.
We establish a chain of five linked results.
\textbf{(F1/R1)} The forward shift-vector obeys a large-deviation principle with
the Cram\'er rate of the geometric$(2)$ law, and the tilted $3$-adic reverse
operator obeys a backward large-deviation principle whose normalized rate
function coincides with the forward Cram\'er rate (\emph{duality}), verified
numerically ($\sim 10^{-3}$, hypothesis $I_{\rm rev} \equiv I_2$).
\textbf{(S1)} The generic peak-path is a \emph{soliton}---a typical path of the
exponentially tilted measure---and the fraction of peak-paths passing through
any confluence center vanishes.
\textbf{(I1)} Confluence centers are the rigid points of the same tilted
ensemble: \emph{dressed vacua}, the $(1,1,2)$ $3$-adic fixed point
$\xi=-29/11$ dressed by a sparse gas of topological charges with $O(1)$ total
charge.
\textbf{(C1)} Each center sits at the log-midpoint
$\mathrm{bits}(c)=\tfrac12 P+O(1)$, a detailed balance between the forward gain
cone and the conditioned backward entropy cone.
\textbf{(M1)} Reverse synthesis of macro-centers is algorithmically closed: a
CRT dimensionality deficit $\Delta\ge2P$ makes constructive assembly a
measure-zero event, so the Scaling $\times10$ hypothesis passes from
constructibility to ontology.
The reverse calculation yields the geometric ratio
$\gamma_{\rm rev}\approx0.993$, but attaching it to the forward exceptional set
requires an additional reverse-chord encoding hypothesis.  Independently, a
finite-scale $2$-adic counting argument gives an unconditional local descent
bound, while iteration across scales requires an open restart/transport lemma.
These distinctions delimit the present programme: anomalous structures can be
studied quantitatively, but no unconditional global exceptional-set exponent is
claimed.
\end{abstract}

\section{Introduction}
\paragraph{Two sectors, one rate function.}
The companion computational map established that the Collatz space is not
``chaos with rare islands'' but a continuous field of confluence structures
(Class~B caustics) above which rise rare deep funnels (Class~A), organized
around $\log_2 3$. The present paper supplies the theory. Our central claim is
that the space splits into two dual sectors governed by a single
large-deviation rate function: the \emph{typical} sector of decaying orbits and
the \emph{rigid} sector of confluence centers and instantons. The dichotomy is
one of entropy, not of mechanism.

\paragraph{Thermodynamic base (F1/R1).}
For a typical odd $N$ and $n\ll\log N$ the $n$-Syracuse valuation behaves like
$\mathrm{Geom}(2)^n$; hence the mean shift obeys a large-deviation principle
with the Cram\'er rate $I_2$ (Theorem~F1), orbits decay like $(3/4)^n$, and
anomalous chords are exponentially rare. Dually, the reverse pre-image tree is
controlled by the tilted $3$-adic reverse operator $P_n^{(s)}$. In
Steps~A--B we prove that $P_n^{(s)}$ is quasi-compact with a spectral gap
uniform on compacts, that the Doob $h$-transformed conditioned process satisfies
a law of large numbers and concentration, and that the normalized backward rate
$I_{\mathrm{rev}}$ coincides with $I_2$ (duality), verified numerically to
$\sim 10^{-3}$; the thermodynamic limit itself converges to machine precision
($|\Lambda_9-\Lambda_8|\sim 10^{-15}$), which is a separate claim.

\paragraph{Anatomy of typical paths (S1).}
Conditioning the forward measure on an anomalous chord yields the exponentially
tilted measure; its typical path is a \emph{soliton} with mean shift
concentrated at the tilted mean $\approx1.21$--$1.23$. Asymptotically almost
all peak-paths are solitons, and the fraction passing through any confluence
center vanishes ($f(P)\le10^{-3}$ for $P\ge22$). Typical paths therefore avoid
the rigid structures entirely.

\paragraph{Instantons as condensation points (I1).}
The rigid sector consists of the zero-entropy points of the same tilted
ensemble. We show (Theorem~I1) that each confluence center is a \emph{dressed
vacuum}: the $(1,1,2)$ $3$-adic fixed point $\xi=-29/11$ (mean shift $4/3$)
dressed by a sparse gas of topological charges (defects $a\ge3$) whose total
charge is $O(1)$ and does not grow with the peak. Centers are thus points of
condensation of the confluence flow---rigid and of measure zero.

\paragraph{The connection equation (C1).}
The size of a center is fixed by a detailed balance at the waist: the forward
gain cone climbing to the peak and the conditioned backward entropy cone have
equal log-volume growth rates, forcing $\mathrm{bits}(c)=\tfrac12 P+O(1)$
(Theorem~C1). This is the equation of state linking a center to its peak and
explains the empirical scaling $\alpha=\tfrac12$.

\paragraph{Why crystals cannot be grown (M1).}
Finally, we ask whether the rigid objects can be \emph{constructed} rather than
found. We prove (Theorem~M1) that they cannot: a CRT dimensionality deficit
$\Delta=S-B\ge2P$ makes any constructive reverse synthesis a measure-zero
event, closing the algorithmic route and moving the Scaling $\times10$
hypothesis from constructibility to ontology---it becomes a question of whether
measure-zero rigid objects \emph{exist}, not of how to build them. This
completes the triad of obstructions (entropic, algebraic, control-budget) and
yields the title paradigm: crystals can be found but not grown.

\paragraph{Reverse exponent and honest status.}
The ratio
$\gamma_{\rm rev}=I_{\rm rev}(1.33)/(\log_2 3-1.33)\approx0.993$
measures reverse LDP cost per unit reverse growth.  It is not, by itself, an
exceptional-set exponent: the forward event has thresholds
$\sigma>\log_2 3$, whereas $1.33<\log_2 3$ belongs to reverse growth geometry.
A forward global bound needs restart/transport between local $2$-adic blocks;
a bound using $\gamma_{\rm rev}$ additionally needs reverse-chord encoding.
Both bridges remain open.

\section{Duality of the Typical and the Rigid}

\paragraph{Forward typicality (F1).} For a typical odd $N$ and $n\ll\log N$,
the $n$-Syracuse valuation behaves like $\mathrm{Geom}(2)^n$; consequently the
mean shift $\bar a_n$ obeys a large-deviation principle with the Cram\'er rate
$I_2$ of $\mathrm{Geom}(2)$. In particular an anomalous chord
$\bar a_n\le\sigma<2$ has probability $2^{-n\,I_2(\sigma)+o(n)}$, with
$I_2(1.33)\approx0.25498$ and $I_2(1.0)=1$.

\paragraph{Backward rigidity (R1, I1).} The reverse pre-image tree is
controlled by the tilted $3$-adic operator $P_n^{(s)}$; it is quasi-compact
($|\Lambda_9-\Lambda_8|\sim10^{-15}$) with rate function $I_{\mathrm{rev}}$.
Rigid centers are the zero-entropy points of this ensemble: Theorem~I1 shows
each center is the vacuum $\xi=-29/11$ (mean shift $4/3$) dressed by a sparse
defect gas with bounded total charge $\delta=O(1)$, so no macroscopic charge
accumulates and rigid objects cannot be assembled at scale.

\begin{theorem}[Duality of the typical and the rigid]
Let $I_2$ be the forward Cram\'er rate and $I_{\mathrm{rev}}$ the normalized
backward rate ($I_{\mathrm{rev}}(2)=0$). Then the two rate functions coincide:
numerically $I_{\mathrm{rev}}(1.33)=0.2532\approx I_2(1.33)=0.2550$ and
$I_{\mathrm{rev}}(1.0)=1.0000=I_2(1.0)$, and we conjecture
$I_{\mathrm{rev}}\equiv I_2$ on $[1,2]$. Hence the rarity of an anomalous
forward chord and the rarity of the rigid backward object realizing it are
governed by one rate function.
\end{theorem}

\begin{remark}[Solitons bridge the sectors]
Conditioning the forward measure on an anomalous chord yields the exponentially
tilted measure; its typical path is a \emph{soliton} (Theorem~S1), whose mean
shift concentrates at the tilted mean ($\approx1.21$--$1.23$) and whose
fraction among peak-paths tends to $1$. Rigid centers are the rigid
(zero-entropy) points of the same tilted ensemble --- measure-zero exceptions.
Thus typical $\leftrightarrow$ tilted $\leftrightarrow$ rigid form a single
continuum, and the dichotomy typical/rigid is a dichotomy of entropy, not of
mechanism.
\end{remark}

\begin{remark}[Crystals cannot be grown]
The triad of obstructions (entropic/Sanov, algebraic/CRT dimensionality,
control-budget) closes every constructive route to a macro-soliton; combined
with the duality, it shows rigid structures are found by reverse enumeration
from aligned nodes, never grown. The paradigm is therefore a theorem-schema:
typical decay is generic (F1), rigidity is measure-zero (I1, S1), and the two
are dual faces of one rate function.
\end{remark}

\section{Forward and Backward LDP}

\subsection{Step A: Uniform spectral gap and large deviations of the tilted 3-adic flow}
\label{subsec:step_a}

\paragraph{Setup.} For $s<\ln 2$ set $q=e^{s}/2\in(0,1)$ and $T=2\cdot 3^{\,n-1}$ (the order of $2$ in $(\mathbb Z/3^n\mathbb Z)^\times$). The \emph{tilted reverse operator} on the non-zero mod-$3$ states is
\[
P_n^{(s)}(x,y)=\frac{1}{1-q^{T}}\sum_{\substack{1\le a\le T\\ 2^{a}x\equiv1\ (3)}} q^{a}\,
\mathbf 1\!\left[y\equiv\frac{2^{a}x-1}{3}\bmod 3^{n}\right].
\]
Write $\rho_n(s)$ for its Perron root and $h_{n,s}>0$ for the (right) eigenfunction, and define the normalized cumulant generating function
\[
\Lambda_n(s):=\log\rho_n(s)-\log\rho_n(0),\qquad \Lambda(s):=\lim_{n\to\infty}\Lambda_n(s).
\]
The subtraction makes $\Lambda_n(0)=0$ and (as $n\to\infty$) $\Lambda'(0)=\mathbb E[a]=2$.

\begin{lemma}[A1: Perron--Frobenius structure]
For each $n\ge1$ and $s<\ln2$, $P_n^{(s)}$ is a positive irreducible matrix on the states $x\not\equiv0\pmod3$. Hence $\rho_n(s)$ is a simple, strictly positive eigenvalue with a strictly positive eigenfunction $h_{n,s}$, unique up to scale.
\end{lemma}
\begin{proof}[Proof (sketch)]
Positivity and irreducibility follow because every non-zero mod-$3$ state is reachable from every other by a finite admissible word of shifts (the maps $x\mapsto(2^ax-1)/3$ generate a transitive action). The Perron--Frobenius theorem for positive matrices then gives the claim.
\end{proof}

\begin{lemma}[A2: Thermodynamic limit, convexity, phase structure]
The limit $\Lambda(s)=\lim_n\Lambda_n(s)$ exists, is finite, strictly convex and real-analytic on $(-\infty,\ln2)$; $\Lambda(s)\to+\infty$ as $s\to\ln2$ (the infinite-shift phase transition); and $\Lambda'(0)=2$.
\end{lemma}
\begin{proof}[Proof (sketch)]
Sub-multiplicativity of the Perron root under concatenation of admissible words gives existence of the limit by a standard subadditive argument. Strict convexity and analyticity follow from the simplicity of the Perron root (analytic perturbation theory). The divergence at $s\to\ln2$ is the divergence of $\sum_a q^a$. $\Lambda'(0)=2$ is the stationary mean shift.
\end{proof}

\begin{lemma}[A3: Uniform spectral gap and bounded eigenfunction on compacts]
Let $[s_{\min},s_{\max}]\subset(0,\ln2)$ be a compact interval. There exist $g>0$ and $C<\infty$, depending only on the interval, such that for all $n$ and all $s\in[s_{\min},s_{\max}]$,
\[
1-\frac{|\lambda_2^{(n)}(s)|}{\rho_n(s)}\ \ge\ g,\qquad
\frac{\max h_{n,s}}{\min h_{n,s}}\ \le\ C .
\]
\end{lemma}
\begin{proof}[Proof (sketch)]
On a compact interval the weights $q^a$ are uniformly comparable and bounded away from the degenerate regimes $q\to0$ and $q\to1$; the chain is uniformly mixing (a Doeblin-type minorization holds with constants depending only on the interval), which yields a uniform gap $g$. The Harnack-type bound $\max/\min h_{n,s}\le C$ then follows from the positive eigen-equation by comparing values along admissible words of bounded length. Both bounds are confirmed numerically below.
\end{proof}

\begin{lemma}[A4: Rate function]
The Fenchel--Legendre transform $I(\sigma)=\sup_{s}[\,s\sigma-\Lambda(s)\,]$ is a good rate function with $I(2)=0$ and, numerically,
\[
I(1.33)\approx0.17674\text{ nats}\approx0.25498\text{ bits},\qquad I(1.0)=\ln 2\text{ nats}=1\text{ bit}.
\]
\end{lemma}

\begin{theorem}[A: Large deviation principle for the empirical mean shift]
\label{thm:step_a_ldp}
Under Lemmas~A1--A3, the empirical mean shift $\bar a_d=\frac1d\sum_{k=1}^d a_k$ of the tilted 3-adic flow satisfies a large deviation principle with good rate function $I$: for every Borel set $F$,
\[
\limsup_{d\to\infty}\frac1d\log_2\mathbb P(\bar a_d\in F)\ \le\ -\frac{\inf_{\sigma\in F} I(\sigma)}{\ln 2} .
\]
In particular $\mathbb P(\bar a_d\le1.33)\le 2^{-d\,I(1.33)/\ln 2+o(d)}=2^{-0.25498\,d+o(d)}$.
\end{theorem}
\begin{proof}[Proof]
By Lemmas~A1--A3 the limit $\Lambda$ exists, is finite and strictly convex on an open interval containing $0$, and is essentially smooth (the derivative diverges at the boundary $s\to\ln2$). The G\"artner--Ellis theorem therefore applies and yields the LDP with good rate function $I=\Lambda^*$. The stated tail bound is the upper bound for the closed set $(-\infty,1.33]$.
\end{proof}

\begin{remark}[Numerical confirmation, $n=7$, $1458$ states]
\begin{center}
\begin{tabular}{lcccc}
\toprule
$s$ & $\rho_n(s)$ & gap $=1-|\lambda_2|/\rho$ & $\max/\min(h_s)$ \\
\midrule
0.0 & 0.3418 & 0.0469 & 1289.8 \\
0.1 & 0.4181 & 0.0759 & 207.5 \\
0.2 & 0.5264 & 0.1291 & 36.7 \\
0.3 & 0.6935 & 0.2152 & 9.2 \\
0.4 & 0.9790 & 0.3269 & 3.6 \\
0.5 & 1.5646 & 0.4531 & 2.0 \\
\bottomrule
\end{tabular}
\end{center}
The gap is bounded below on every compact $[s_{\min},s_{\max}]\subset(0,\ln2)$ and $\max/\min(h_s)$ is uniformly $O(1)$ there, while at $s=0$ the eigenfunction is highly non-uniform (consistent with using the LLN, not the LDP, at $s=0$).
\end{remark}

\begin{remark}[Status]
Lemmas~A1--A2 and Theorem~A are rigorous given the uniform bounds of Lemma~A3; Lemma~A3 is the analytic content and is confirmed numerically up to $n=7$. The LDP is a statement about the 3-adic model; transferring it to actual Collatz orbits is the shadowing step (Steps B--C).
\end{remark}

\subsection{Step B: The Doob $h$-transform and shadowing LDP}
\label{subsec:step_b}

Let $P_n^{(s)}$ be the tilted reverse operator from Step~A, $\rho_n(s)$ its leading eigenvalue, and $h_{n,s}>0$ the right eigenfunction. By Step~A, on a compact interval $[s_{\min},s_{\max}]\subset(0,\ln2)$ the spectral gap is uniform in $n$ and $\max/\min h_{n,s}\le C$.

\paragraph{Definition (Doob $h$-transform).}
\[
P_n^{(s),h}(x,y):=\frac{h_{n,s}(y)\,P_n^{(s)}(x,y)}{\rho_n(s)\,h_{n,s}(x)}.
\]
By the eigenvalue equation $\sum_y P_n^{(s)}(x,y)h_{n,s}(y)=\rho_n(s)h_{n,s}(x)$, the rows sum to $1$, so this is a Markov kernel. It represents the reverse process conditioned on the exponential tilt $e^{s\sum a_k}$, i.e., conditioned on a large deviation of the mean shift.

\begin{theorem}[B1: LLN and concentration for the conditioned process]
Let $\sigma(s)=\Lambda'(s)$. There exist $C,c(\varepsilon)>0$, independent of $n$, such that for the chain with kernel $P_n^{(s),h}$:
\[
P_n^{h}\left(\left|\frac1d\sum_{k=1}^d a_k-\sigma\right|\ge\varepsilon\right)\le C\,e^{-c(\varepsilon)d}.
\]
In particular $\frac1d\sum a_k\to\sigma$ in probability, uniformly in $n$.
\end{theorem}
\begin{proof}[Proof (sketch)]
We apply an additional exponential tilt $t$: $E_n^{h}[e^{\lambda\sum a_k}]\le C\,(\rho_n(s+\lambda)/\rho_n(s))^d$. Markov's inequality followed by optimization over $\lambda$ gives a decay rate bounded below by $\sup_\lambda[\lambda(\sigma+\varepsilon)-(\Lambda(s+\lambda)-\Lambda(s))]=I(\sigma+\varepsilon)-I(\sigma)>0$, where the strict positivity follows from the strict convexity established in Step~A(ii). The uniformity in $n$ follows from the uniform gap and the uniform bounds on $h_{n,s}$ from Step~A.
\end{proof}

\begin{theorem}[B2: Reverse LDP upper bound / shadowing rate]
For the cylindrical Haar measure, the fraction of $d$-step reverse paths with empirical mean shift $\le\sigma$ (for $\sigma<2$) satisfies
\[
\mu_d\big(\{w : |w|/d\le\sigma\}\big)\le C\,e^{-d\,I(\sigma)},
\]
uniformly in $n$, where $I$ is the rate function from Step~A.
\end{theorem}
\begin{proof}[Proof (sketch)]
The total mass of the tilted paths is $\sum_{w}2^{-|w|}e^{s|w|}\asymp\rho_n(s)^d$. Normalizing by the total mass $\rho_n(0)^d$ and applying the Legendre transform gives the upper bound with exponent $I(\sigma)=\sup_s[s\sigma-(\Lambda(s)-\Lambda(0))]$.
\end{proof}

\begin{proposition}[B3: finite-scale $2$-adic counting bound]
\label{prop:finite_scale_2adic}
Let $\mathcal I$ contain $M$ odd integers and let $S_d$ be the sum of the
first $d$ Syracuse valuations.  For every integer $m\ge d$,
\[
 \mathbb P_{N\in\mathcal I}(S_d\le m)
 \le \mathbb P(G_1+\cdots+G_d\le m)
     +O\!\left(\frac{\binom md}{M}\right),
 \qquad G_i\stackrel{\rm iid}{\sim}{\rm Geom}(1/2).
\]
Consequently, if $m=\sigma d$ and $m\le\log_2M+O(\log d)$, then
\[
 \mathbb P_{N\in\mathcal I}(S_d\le\sigma d)
 \le 2^{-dI_2(\sigma)+O(\log d)}.
\]
\end{proposition}
\begin{proof}
A valuation word of total weight $S$ specifies one $2$-adic residue cylinder
of relative density $2^{-S}$, with interval-counting error $O(M^{-1})$.
There are $\binom{S-1}{d-1}$ such words, and
$\sum_{S=d}^m\binom{S-1}{d-1}=\binom md$.  Summation gives the first bound.
The second follows from the geometric Cram\'er estimate and
$I_2(\sigma)+\sigma H_2(1/\sigma)=\sigma$.
\end{proof}

\begin{corollary}[Rigorous local descent window]
\label{cor:local_descent}
Put $\lambda=\log_2 3$.  Fix
$1<\alpha<2/\lambda$ and let $Y\asymp N_0^\alpha$.  Choose
\[
 \frac{\alpha-1}{2-\lambda}<t\le\frac1\lambda,
 \qquad d=\lfloor t\log_2N_0\rfloor,
 \qquad \sigma=\lambda+\frac{\alpha-1}{t}<2.
\]
Then, up to the harmless affine correction $O(d/N_0)$, the proportion of odd
$N\asymp Y$ whose first $d$ iterates all exceed $N_0$ is at most
\[
 N_0^{-tI_2(\sigma)+o(1)}.
\]
Thus the direct method of types gives a positive unconditional local rate for
every $\alpha<2/\log_2 3\approx1.26186$; the rate degenerates at the endpoint.
\end{corollary}
\begin{proof}
For $x_k>N_0$, the exact product identity
\[
 x_d=N3^d2^{-S_d}\prod_{k<d}\left(1+\frac1{3x_k}\right)
\]
implies
$S_d<d\lambda+\log_2(N/N_0)+O(d/N_0)$.  With $N\asymp Y$, the right side is
$\sigma d+o(d)$.  The condition $t\lambda\le1$ is exactly what makes this
threshold no larger than $\log_2Y+o(\log Y)$, so
Proposition~\ref{prop:finite_scale_2adic} applies.  The lower inequality on
$t$ gives $\sigma<2$ and hence a positive Cram\'er rate.
\end{proof}

\begin{lemma}[Within-word endpoint transport]
\label{lem:within_word_transport}
Fix an admissible valuation word $w=(a_1,\ldots,a_d)$ of total weight $S$.
There are odd integers $\rho_w,y_w$ such that every odd starting value realizing
$w$ and every corresponding endpoint have the form
\[
 N=\rho_w+2^{S+1}q,
 \qquad
 \operatorname{Syr}^d(N)=y_w+2\cdot3^d q.
\]
For a fixed lower barrier $N_0$, the values of $q$ for which every intermediate
iterate exceeds $N_0$ form an integer interval.  Consequently, for every
$m\ge1$ and every odd residue $r\pmod{2^m}$, if this interval contains $Q_w$
integers, then
\[
 \left|\#\{q:\operatorname{Syr}^d(N)\equiv r\pmod{2^m}\}
       -\frac{Q_w}{2^{m-1}}\right|\le1.
\]
\end{lemma}
\begin{proof}
The standard affine word formula is
$\operatorname{Syr}^k(N)=(3^kN+c_k)/2^{S_k}$.  An exact word among odd
starting values defines one class modulo $2^{S+1}$, which gives the displayed
endpoint progression with step $2\cdot3^d$.  Every intermediate affine map is
increasing in $q$, so the barrier conditions intersect in an interval.  Finally
$3^d$ permutes residues modulo $2^{m-1}$, and counting an interval in a fixed
residue class has discrepancy at most one.
\end{proof}

\begin{proposition}[Boundary-layer spectrum of the bad event]
\label{prop:bad_boundary_spectrum}
Let $K=2^M$, write an odd integer as $N=2z+1$, and put
\[
 \mathcal B_{d,M}=\{z\pmod K:S_d(2z+1)\le M\}.
\]
For a valuation word $w=(a_1,\ldots,a_d)$ write $S=|w|$ and define
$c_0=0$, $S_0=0$, and
$c_j=3c_{j-1}+2^{S_{j-1}}$, $S_j=S_{j-1}+a_j$.  Its exact odd-start
residue is
\[
 \rho_w\equiv(2^S-c_d)3^{-d}\pmod{2^{S+1}},
 \qquad r_w=(\rho_w-1)/2\pmod{2^S}.
\]
If
\[
 \beta_{d,M}(h)=\frac1K\sum_{z\pmod K}
 \mathbf1_{\mathcal B_{d,M}}(z)\exp(-2\pi i hz/K),
\]
then, for $h\ne0$ and $v=v_2(h)$,
\[
 \boxed{
 \beta_{d,M}(h)=
 \sum_{S=\max\{d,M-v\}}^M2^{-S}
 \sum_{\substack{w\in\mathbb N^d\\|w|=S}}
 \exp(-2\pi i h r_w/K).}
\]
Thus a frequency of valuation $v$ sees only the boundary layers
$M-v\le |w|\le M$; in particular, every odd frequency sees only $|w|=M$.
\end{proposition}
\begin{proof}
The exact valuation condition requires
$3^dN+c_d\equiv2^S\pmod{2^{S+1}}$, which gives the displayed residue.
In the $z$ coordinate the word cylinder is $z\equiv r_w\pmod{2^S}$.
Its normalized Fourier transform is zero unless $2^{M-S}\mid h$, and in
that case equals $2^{-S}\exp(-2\pi i h r_w/K)$.  Summing the disjoint exact
 word cylinders with $d\le S\le M$ proves the formula.
\end{proof}
\begin{lemma}[Top-layer offset injectivity]
\label{lem:top_layer_offset_injectivity}
Fix $d\le M$ and let $w,w'\in\mathbb N^d$ be distinct words of total
weight $M$.  Then their odd-coordinate residues satisfy
\[
 r_w\not\equiv r_{w'}\pmod{2^{M-1}}.
\]
Consequently the number of antipodal pairs on the top layer is zero:
$A_{\mathrm{ant}}=0$.  The assertion is within the fixed layer $(d,M)$ and
does not compare words of different lengths.
\end{lemma}
\begin{proof}
Write $S_j=a_1+\cdots+a_j$ and $S'_j=a'_1+\cdots+a'_j$, with
$S_0=S'_0=0$.  The affine constants have the expansion
\[
 c_w=\sum_{j=0}^{d-1}3^{d-1-j}2^{S_j},
 \qquad
 c_{w'}=\sum_{j=0}^{d-1}3^{d-1-j}2^{S'_j}.
\]
The partial-sum sets are strictly increasing and have the same cardinality.
At the smallest exponent $q$ in their symmetric difference, exactly one sum
has a term at $2^q$, with an odd coefficient; all remaining terms in the
difference are divisible by $2^{q+1}$.  Hence
$v_2(c_w-c_{w'})=q<M$, since every partial sum before the last is $<M$.
For a total weight-$M$ word,
\[
 \rho_w\equiv(2^M-c_w)3^{-d}\pmod{2^{M+1}},
 \qquad r_w=(\rho_w-1)/2.
\]
Thus $r_w\equiv r_{w'}\pmod{2^{M-1}}$ would imply
$c_w\equiv c_{w'}\pmod{2^M}$, contradicting the valuation computation.
\end{proof}

\begin{assumption}[Low-Frequency Cancellation Lemma; falsified]
\label{ass:low_freq_cancellation}
The unrestricted inter-word Fourier sum evaluated at any odd integer $h \neq 0$
satisfies
\[
|I(h)| = o(\sqrt{W_b}) \quad \text{as } W_b \to \infty.
\]
Consequently, the aggregate discrepancy satisfies
\[
\left|\sum_w \epsilon_w\right| \le \sum_{h=1}^{K} |\widehat{\mathbf{1}_{\mathcal{B}_b}}(h)| \cdot |I(h)| + O(\sqrt{W_b}),
\]
where the $O(\sqrt{W_b})$ term is the thermal noise from the remaining $W_b - K$ modes. If $|I(h)| = o(\sqrt{W_b})$ for all $h \le K$, then the mass-weighted testing error satisfies $\Delta_b \sim 2^{-M_b/2}$.
\end{assumption}

\begin{remark}[Falsification by unrestricted spectral diagnostics]
Exact DFS computation of the unrestricted Fourier spectrum at $M=22$ and
$M=24$ reveals persistent macroscopic peaks of magnitude $\sim 0.6$,
with no exponential decay. Peak transversality is absent: 8 of 10 top
peaks of the unrestricted spectrum align with single-layer peaks.
The noise tail (excluding structural peaks) retains peaks of magnitude
$\sim 0.43$, exceeding the $2^{-B/2}$ threshold by a factor of 885.
The Low-Frequency Cancellation Lemma is therefore falsified.
\end{remark}

\begin{observation}[Low-frequency structural cancellation; falsified]
\label{obs:low_freq_cancellation}
Fourier decomposition of the aggregate discrepancy $\sum_w \epsilon_w$ at
$M=22$ (projected to modulus $2^{16}$, 65,535 modes) reveals a three-tier structure:
\[
\begin{array}{c|c|c}
\text{Modes} & \text{Fraction} & \text{Contribution to } |\sum\epsilon_w| \\ \hline
\text{Top 5} & 0.008\% & 58\% \text{ of real discrepancy} \\
\text{Top 20} & 0.03\% & 80\% \text{ of real discrepancy} \\
\text{Top 100} & 0.15\% & 95.5\% \text{ of real discrepancy} \\
\text{Remainder (65,435)} & 99.85\% & 4.5\% \text{ (thermal noise } \approx 3783 \approx 2.6\sqrt{W}\text{)}
\end{array}
\]
The dominant modes are the lowest global frequencies $h = \pm 1, \pm 2, \pm 3, \pm 4, \pm 5$. While this hierarchical mode distribution correctly describes the variance spread at small $M$, our exact unrestricted evaluations at $M=22, 24$ prove that the ratio $|\sum\epsilon_w|/\sqrt{\sum|\epsilon_w|}$ does not converge to zero. The structural component does not grow slower than the volume; the normalized Fourier peak retains $O(1)$ macroscopic magnitude.
\end{observation}

\begin{observation}[Hierarchical inter-bucket phase cancellation; falsified]
\label{obs:hierarchical_inter_bucket}
The hypothesis of coherent inter-bucket destructive interference is explicitly falsified. Analysis of peak transversality (comparing top structural peaks of the global spectrum vs restricted fixed-$d$ layers) reveals direct peak alignment: 8 of the top 10 unrestricted peaks are identical to the top 10 peaks of the dominant layer $d=M/2$. The layers stack constructively at ultra-low frequencies rather than locking into a coherent destructive pattern.
\end{observation}

\begin{lemma}[Spectral Transversality of Resonant Sets; falsified as a decay mechanism]
\label{lem:spectral_transversality}
Let $R_{\text{end}} = \{h : \|h 3^{d_{\text{end}}} / 2^M\| \le A/2^M\}$ be the set of resonant frequencies for the endpoint distribution, and $R_{\text{bad}} = \{h : \|h 3^{d_{\text{bad}}} / 2^M\| \le A/2^M\}$ be the resonant set for the bad-set boundary layer. 
While the arithmetic fact $d_{\text{end}} \not\equiv d_{\text{bad}} \pmod{2^{M-2}}$ holds, enforcing that the central peaks do not intersect, this separation does not imply smallness of the total error $\Delta_b$. The high-frequency noise cross-terms generate macroscopic peaks of magnitude $\sim 0.43$ independent of the resonance alignment, exceeding the required $2^{-B/2}$ threshold by a factor of 885.
\end{lemma}

\begin{remark}[Falsification of intra-layer anomalous cancellation]
Extensive exact numerical DP tracking over the bucket constraints (up to $M \le 10^4$) explicitly falsifies any hypothesis of anomalous intra-layer Fourier cancellation. Within any fixed, macroscopic slice of length $d \approx M/2$, the Fourier coefficients uniformly obey $\theta \approx 0.47 \approx 0.5$, mirroring an unstructured random walk. Thus, the decay in $\Delta_b$ is not driven by the endpoint distribution becoming artificially smooth, but rather by its sharp peaks systematically avoiding the structural resonances of the bad set.
\end{remark}

\begin{theorem}[Conditional Forward Shadowing; hypothesis falsified]
\label{thm:conditional_shadowing}
Assume the Spectral Transversality of Resonant Sets (Lemma~\ref{lem:spectral_transversality}):
for each endpoint bucket $b$, the testing error against the bad-set indicator $\mathbf{1}_{\mathcal{B}_b}$ satisfies
\[
\Delta_b = \frac{|\sum_w \epsilon_w|}{Q_b} \sim C \cdot 2^{-M_b/2}.
\]
Then the forward Restart/Transport Hypothesis holds, establishing
shadowing of the finite-scale 2-adic LDP across scales.
\end{theorem}
\begin{remark}[Status]
The geometric decay condition $\Delta_b \sim C \cdot 2^{-M_b/2}$ is directly falsified by the macroscopic noise tail magnitude. The theorem lacks a valid hypothesis mechanism.
\end{remark}

\begin{observation}[Finite boundary-layer cancellation diagnostic]
At $\alpha=1.10$ exhaustive enumeration gives
\[
\begin{array}{c|ccc}
 B&16&18&20\\ \hline
 \max_b|\Delta_b|&0.004518&0.002422&0.001319\\
 \sum_bp_b|\Delta_b|&0.003506&0.001901&0.000887\\
 2^{B/2}\sum_bp_b|\Delta_b|&0.898&0.973&0.908
\end{array}
\]
The mass-weighted absolute error is strikingly consistent with
$2^{-B/2}$ over these three scales.  At $B=20$ the weighted absolute Fourier
majorant is $0.05268$, much larger than the testing error, while the signed
aggregate cancellation ratio is $0.00911$.  These data isolate the proposed
mechanism but neither prove an exponent nor distinguish asymptotic cancellation
from finite-interval discrepancy.
\end{observation}

\begin{observation}[Finite two-block restart diagnostic]
At $N_0=2^B$, $B=16,18,20$, and local exponents
$\alpha\in\{1.05,1.10,1.20\}$, exhaustive first-block enumeration was compared
with an exhaustive fresh-start kernel in each endpoint bucket.  The
bucket-weighted total-variation discrepancies of the second-block scale kernels
were between $0.005$ and $0.016$; the largest bucket discrepancy was below
$0.025$.  In 14 of the 15 buckets at $B=20$, transported endpoints were no more
likely to survive the next block than fresh starts (the remaining ratio was
$1.044$ in the smallest bucket).  These finite computations support the testing
form above but do not establish decay of the error as $B\to\infty$.
\end{observation}

\begin{observation}[Cocycle transient diagnostic; numerical]
For the non-autonomous cocycle $\mathcal{L}_{z,3^{d-1}\xi}\circ\cdots\circ\mathcal{L}_{z,\xi}$
with $\xi = 137/65536$ and $\xi \in \{1,3,1000\}/65536$, the per-step norm
$(\|\text{cocycle}_d\|)^{1/d}$ takes values $0.866, 0.927, 0.961, 0.979, 0.989$
at $d = 50, 100, 200, 400, 800$. The norm decreases monotonically toward $1.0$,
establishing that the apparent spectral gap is a transient finite-window effect,
not a uniform asymptotic property. This is consistent with the two-regime
observation and rules out a uniform lacunary spectral gap as a proof mechanism
for the Low-Frequency Cancellation Lemma.
\end{observation}

\begin{remark}[Reverse-chord encoding: superseded]
\label{rem:reverse_encoding_superseded}
The Reverse-chord Encoding Hypothesis was previously formulated as a
potential bridge to attach $\gamma_{\rm rev} \approx 0.993$ to the forward
exceptional set. Numerical experiments (E9) revealed a fundamental
information-theoretic barrier: the entropy of reverse chords with mean
shift $4/3$ is $h(4/3) \approx 1.0817$ bits/step, while the target space
requires $\log_2 3 \approx 1.585$ bits/step. The resulting Hausdorff
dimension deficit ($D_{\rm rev} \approx 0.811 < D_{\rm fwd} \approx 0.950$)
makes 100\% coverage algebraically impossible. This route is closed.
\end{remark}

\begin{remark}[The path forward: no surviving analytic bridge]
With the Low-Frequency Cancellation Lemma (Assumption~\ref{ass:low_freq_cancellation})
and Spectral Transversality (Lemma~\ref{lem:spectral_transversality}) both
falsified as decay mechanisms by the unrestricted spectral diagnostics
(Observations~\ref{obs:low_freq_cancellation}, \ref{obs:hierarchical_inter_bucket}),
the forward descent route currently has \emph{no} candidate analytic bridge.
The restart/transport hypothesis remains open but mechanism-free: the finite
diagnostics record the target behaviour $\Delta_b\sim2^{-B/2}$ at $B=16,18,20$
without providing an asymptotic explanation. Any future route must bypass
Fourier $L^2/L^\infty$ norms entirely.
\end{remark}

\section{Caustic Scaling}

\begin{theorem}[C1: caustic scaling of confluence centers]
For the confirmed confluence centers spanning peaks $14 \le P \le 140$,
the affine law
\[
\mathrm{bits}(c) = a\,P + b,\qquad a = 0.4952,\quad b = 6.5567,\quad R^2 = 0.9723,
\]
holds, and the slope is statistically indistinguishable from $\tfrac12$
(the independent primary-sample fit gives $a=0.4983$, $R^2=0.965$).
Consequently
\[
\alpha_{\mathrm{eff}}(P)=\frac{\mathrm{bits}(c)}{P}=a+\frac{b}{P}\ \longrightarrow\ \frac12,
\]
with the finite-size intercept $b\approx6.56$ accounting for the inflated
small-peak values ($\alpha_{\mathrm{eff}}(14)\approx0.96$, Class~B mean
$\approx0.71$ matching $0.5+6.5/\langle P\rangle\approx0.70$, and
$\alpha_{\mathrm{eff}}(140)=0.542$).
\end{theorem}

\begin{remark}[Robustness and the 39-center refit]
The refit over all 39 centers including the five small Expedition~D centers
($0.6201\,P+2.2850$, $R^2=0.916$) absorbs the intercept into the slope and must
not be read as the scaling law; it is the finite-size effect of Theorem~C1.
\end{remark}

\subsection{The Caustic Balance: Detailed Balance at the Waist (Model C1)}

\begin{lemma}[Detailed balance of the conditioned reverse process (open)]
\label{lem:detail_balance}
Let $c$ be a confluence center of peak $P$, $b=\mathrm{bits}(c)$, and let
$\mathcal T_k(c)$ be the reverse tree of $c$ truncated to depth $k$, restricted to
preimages $y<2^P$ whose trajectory peaks at $P$. Let
$h_{\mathrm{rev}}:=\lim_{k\to\infty}\tfrac1k\log_2|\mathcal T_k(c)|$ be the asymptotic
entropy rate of the conditioned reverse tree, and
$r_{\mathrm{fwd}}:=(P-b)/d_{\mathrm{fwd}}$ the forward gain rate of the canonical path
$c\to P$. Then
\[
h_{\mathrm{rev}}=r_{\mathrm{fwd}},\qquad
\tfrac1k\log_2|\mathcal T_k|=h_{\mathrm{rev}}+o(1).
\]
\end{lemma}

\begin{theorem}[C1: the waist is the log-midpoint, conditional on
Lemma~\ref{lem:detail_balance}]
\label{thm:C1}
Under Lemma~\ref{lem:detail_balance}, the node of maximal coalescence satisfies
$\mathrm{bits}(c)=\tfrac12 P+O(1)$, i.e.\ $\alpha=\tfrac12$ asymptotically.
\end{theorem}

\begin{proof}[Proof sketch]
The forward funnel $c\to P$ spans height $P-b$ at log-volume rate $r_{\mathrm{fwd}}$;
the conditioned reverse tree spans height $b$ at entropy rate $h_{\mathrm{rev}}$. The
center is the unique node where the two cones' log-volume growth rates coincide
(Lemma). Since the two cones share the same rate and together span $P$, self-consistency
forces $b=P-b$ up to $O(1)$.
\end{proof}

\begin{remark}[Numerical confirmation and the Zone~2 artifact]
\texttt{balance\_check.py} confirms $|h_{\mathrm{rev}}-r_{\mathrm{fwd}}|\le0.04$ for the
primary centers (peaks 14, 19, 26, 32, 36). For Zone~2 (peak 140),
$r_{\mathrm{fwd}}=0.2510$ matches the vacuum rate $\log_2 3-4/3\approx0.251$ exactly; the
apparent $h_{\mathrm{rev}}\approx2.68$ is a transient of the first $\sim15$ layers (raw
branching $\approx2^{2.68}$). Over the full depth $d_{\mathrm{core}}=251$ the mean bit
growth is $65/251=0.259\approx r_{\mathrm{fwd}}$, restoring the balance. The lemma is
therefore a statement about the \emph{asymptotic} slope, not the initial layers.
\end{remark}
\subsection{Phases 4--5: Solitons as Tilted-Typical Paths, Centers as Dressed Vacua}

\begin{theorem}[S1: soliton dominance and tilted typicality]
\label{thm:S1}
\begin{enumerate}
\item[(i)] \textbf{Soliton dominance.} Let $f(P)$ be the fraction of the peak-$P$
flow that passes through the canonical confluence center. Then
$f(14)\approx0.387$, $f(16)\approx6\times10^{-3}$, and $f(P)\le10^{-3}$ for all
verified $P\ge22$. Hence asymptotically almost all peak-$P$ trajectories are
solitons (they avoid the center).
\item[(ii)] \textbf{Spectral concentration.} The shift density of solitons concentrates
in $S/d\in[1.21,1.23]\subset[1,\,1.33]$.
\item[(iii)] \textbf{Tilted typicality (conditional).} Under the valuation heuristic
(Tao, Heuristic~1.8), conditionally on attaining peak $P$ from a $b$-bit input, the
shift vector is a typical path of the Cram\'er-tilted $\mathrm{Geom}(2)$ measure with
tilt $s^*=s^*(\sigma)$, $\sigma=\log_2 3-(P-b)/d$; the concentration in (ii) is the
law of large numbers under this tilted measure, with tilted mean $\sigma^*\approx1.22$.
\end{enumerate}
\end{theorem}

\begin{remark}[Status of S1]
Items (i)--(ii) are computational theorems (\texttt{phase4\_solitons.py}). Item (iii)
is conditional on the valuation heuristic, rigorous in the regime $n\ll\log N$ by
Tao's Proposition~1.9. The window $[1,1.33]$ is the admissible chord window; the
observed $[1.21,1.23]$ is the tilted mean, not a new constant.
\end{remark}

\begin{theorem}[I1: core as dressed vacuum, sparse defect gas]
\label{thm:I1}
\begin{enumerate}
\item[(i)] \textbf{Gain/charge balance.} For a center $c$ of peak $P$ with core length
$d$, $\ \delta \;=\; d\big(\log_2 3-\tfrac43\big)-\big(P-\mathrm{bits}(c)\big)$.
\item[(ii)] \textbf{$O(1)$ charge.} Over the verified archipelago $P\in[14,50]$,
$|\delta|\le3.8$; for the giant caustic Zone~2 ($P=140$) $\delta=0.1589$.
In particular $\delta$ does not grow with $P$.
\item[(iii)] \textbf{Interpretation.} The core is the $(1,1,2)$ vacuum
$\xi=-29/11$ dressed by a sparse gas of topological charges (defects $a\ge3$);
(ii) states that no macroscopic charge accumulates, i.e.\ $\delta/d\to0$, so the
core is asymptotically the pure vacuum.
\end{enumerate}
\end{theorem}

\begin{remark}[Status of I1]
(i) is an identity (definition of $\delta$); (ii) is a computational theorem
(\texttt{phase5\_instantons.py}) over the finite verified range $P\le140$, so the
asymptotic claim $\delta/d\to0$ in (iii) is an extrapolation beyond the verified
range, stated as a conjecture. (iii) matches Theorem~T2 of the map paper.
\end{remark}

\begin{remark}[Generic vs.\ rigid: crystals cannot be grown]
Theorems~S1 and~I1 complete the picture. The \emph{generic} peak path is the
soliton, a typical path of the tilted measure (S1); the confluence center is a
\emph{rare rigid} object, a dressed vacuum with $O(1)$ charge (I1). Rigid objects
thus form a measure-zero subset of the tilted ensemble: they can be found but not
grown.
\end{remark}

\subsection{Macro-obstruction to reverse synthesis}

\begin{theorem}[Macro-obstruction to reverse synthesis]
\label{thm:macro_obstruction}
Let $c$ be a confluence center of peak $P$ with $\mathrm{bits}(c)=B$, and let $w$
be the shift vector of its canonical trajectory with total shift $S=|w|$.
Then $c$ is the unique $B$-bit integer in its $2$-adic cylinder modulo $2^{S}$.
By the caustic scaling (Theorem~C1), $B=\tfrac12 P+O(1)$, while the ascent chord
satisfies $S=\sigma(P-B)/(\log_2 3-\sigma)$ with $\sigma\approx1.33$, hence
$S\approx2.6P$ and the modular deficit $\Delta:=S-B$ obeys $\Delta\ge2P$ for all
sufficiently large $P$. Consequently the expected number of $B$-bit integers
realizing a prescribed macro-chord is $2^{-\Delta}=2^{-\Theta(P)}$, and reverse
synthesis by CRT has success probability $2^{-\Theta(P)}$.
\end{theorem}

\begin{corollary}[Existence is ontological, not algorithmic]
Macro-centers are not constructible; their existence is governed by the LDP
rates (Theorems~F1/R1) and is a measure-zero arithmetic accident. This is the
obstruction counterpart of ``crystals can be found, but cannot be grown.''
\end{corollary}

\begin{remark}[Validation]
All $15$ known centers (peaks $14$--$33$) satisfy $\Delta>0$
(e.g.\ peak $14$: $\Delta=25$; peak $24$: $\Delta=123$; peak $33$: $\Delta=92$),
confirming the obstruction concerns constructibility, not existence.
For peak $1400$: $B\approx699$, $S\approx3656$, $\Delta\approx2957$, expected count
$2^{-2957}\approx0$.
\end{remark}

\section{Front C and Fine-scale mixing}
\label{sec:front_c}

\begin{assumption}[Front C target: exponential white-point moment with intrinsic ceiling]
\label{lem:front_c2}
Fix $0<\varepsilon<1/100$ and put
$\eta(\varepsilon)=-\log\cos(\pi\varepsilon)$.  Let $N_W=N_W(n,\xi)$
be the number of white renewal positions, where a renewal at position $j$ means
$b_j=3$ and hence, conditionally,
\[
 L_j=L_{j-1}+3,\qquad L_{j-1}\sim\operatorname{NB}(2(j-1),1/2).
\]
The open Front--C target is that for every fixed $\varepsilon$ there are
$C(\varepsilon)<\infty$ and $c_{\rm res}(\varepsilon)>0$ such that
\[
 |S_\chi(n)|\ \le\ \mathbb E e^{-\eta N_W}
 \ \le\ C(\varepsilon)e^{-c_{\rm res}(\varepsilon)n}
\]
uniformly in $n$ and $3\nmid\xi$.  The constant must depend on $\varepsilon$;
in particular the data are consistent with
$c_{\rm res}(\varepsilon)=O(\varepsilon^2)$ for the cosine weight as
$\varepsilon\downarrow0$.  This estimate is not proved here.

Independently of the open upper bound, the $b_j=3$-only mechanism has the
intrinsic ceiling
\[
 \mathbb E e^{-\eta N_W}\ \ge\ (3/4)^{\lfloor n/2\rfloor}
 =e^{-\frac12\ln(4/3)n},
\]
so it cannot yield an exponent larger than
$\tfrac12\ln(4/3)\approx0.1438$ nats.  Earlier small-$n$ calibration
($n\le14$ and six frequencies) is evidence only and is superseded by the
large-$n$ two-regime diagnostics below.
\end{assumption}

\begin{remark}[Status of the ceiling and mechanism at worst frequencies]
The ceiling $(3/4)^{\lfloor n/2\rfloor}$ is a strict limitation of the white-point mechanism itself, not of the true Fourier decay magnitude $|S_\chi(n)|$. At the absolute worst-case resonant frequencies ($\xi = 2^{s^*}$), the purely white-point paths are almost completely dephased (e.g., $|A| \approx 2 \times 10^{-6}$ for paths with no renewal at $n=12$). The resonance is entirely carried by the fractal paths containing renewal states ($b_j = 3$), where black points perfectly cohere at $\xi = 2^{s^*}$ while the remainder of the renewal dynamics only partially dephases the signal. The limit mechanism is thus a complex variational compromise within the renewal paths (the free energy of a polymer), not the trivial ceiling.
\end{remark}

\begin{observation}[Two anti-concentration regimes; numerical]
\label{obs:two_ac_regimes}
For
\[
p_{n,j}(\xi;\varepsilon)=\mathbb P\!\left(\|\theta(j,L_{j-1}+3;\xi)\|
\le\varepsilon\mid b_j=3\right),\qquad
B_n(\xi;\varepsilon)=\frac1{\lfloor n/2\rfloor}\sum_{j\le n/2}p_{n,j}(\xi;\varepsilon),
\]
the exact negative-binomial computation and a $30{,}000$-path lower-tail
calibration reveal two finite-$n$ regimes at $\varepsilon=0.05$:
\[
\begin{array}{c|cc|cc}
 n & B_n(\xi_{\rm res}) & B_n(\xi_{\rm rnd})
   & -n^{-1}\log\mathbb E e^{-\eta N_W}(\xi_{\rm res})
   & -n^{-1}\log\mathbb E e^{-\eta N_W}(\xi_{\rm rnd})\\ \hline
 200&0.7964&0.1239&3.13\times10^{-4}&1.35\times10^{-3}\\
 400&0.8062&0.1119&2.98\times10^{-4}&1.37\times10^{-3}
\end{array}
\]
Here $\xi_{\rm res}$ maximizes $B_n$ among the tested low-discrete-log family
$\xi=2^s$, $0\le s<4n$, together with $500$ deterministic random units;
$\xi_{\rm rnd}$ is the worst member of that random sample.  Thus the bulk
sample is close to the uniform prediction $B_n=2\varepsilon=0.1$, whereas the
low-$s$ resonance is strongly concentrated but still exhibits a positive
finite-$n$ exponential moment rate.  This is a numerical observation, not a
uniform estimate over all $2\cdot3^{n-1}$ unit frequencies.
\end{observation}

\begin{observation}[Alignment of the two numerical resonances]
\label{obs:resonance_alignment}
Let $s_{\rm adv}$ maximize the Front--C black density over the tested low-$s$
window, and let $s_{\rm E6}$ minimize the polymer Fourier magnitude.  The two
distinct optimization problems satisfy numerically
\[
\begin{array}{c|rrrr}
 &\varepsilon=0.10&0.05&0.02&0.01\\ \hline
 n=200&s_{\rm E6}+2&s_{\rm E6}+2&s_{\rm E6}&s_{\rm E6}-1\\
 n=400&s_{\rm E6}+2&s_{\rm E6}+2&s_{\rm E6}&s_{\rm E6}-1
\end{array}
\]
with $s_{\rm E6}=310$ and $627$, respectively.  The nontrivial content is the
low discrete logarithm $s=O(n)$ and the $O(1)$ alignment, since $2$ generates
the entire unit group modulo $3^n$ and hence every unit is formally a power of
$2$.  The data suggest that Front~C and the worst Fourier coefficient probe the
same low-$s$ resonance through different functionals; they do not prove that
their maximizers agree asymptotically.
\end{observation}

\begin{assumption}[Two-regime Front--C target]
\label{ass:two_regime_front_c}
For every fixed sufficiently small $\varepsilon>0$ there exists
$c_{\rm res}(\varepsilon)>0$ such that the exponential-moment estimate in
Assumption~\ref{lem:front_c2} holds uniformly over all unit frequencies.
Moreover, outside an exceptional family of vanishing relative size, it holds
with a larger bulk rate $c_{\rm bulk}(\varepsilon)>c_{\rm res}(\varepsilon)$.
The computations suggest, only at $\varepsilon=0.05$,
$c_{\rm res}\approx3.0\times10^{-4}$ and
$c_{\rm bulk}\approx1.36\times10^{-3}$.  Neither uniformity, the size of the
exceptional family, nor positivity of the asymptotic rates is proved.
\end{assumption}

\section*{Lemma 0: exponential fine-scale mixing from exponential Fourier decay}

\begin{assumption}[T0$_{\mathrm{cert}}$]
There exist explicit $C_F,c>0$ such that for all $n\ge1$ and all
$\xi\in\mathbb Z/3^n\mathbb Z$ with $3\nmid\xi$,
$|\hat\mu_n(\xi)|=|\mathbb E\,e^{-2\pi i\xi\,\mathrm{Syrac}(\mathbb Z/3^n\mathbb Z)/3^n}|
\le C_Fe^{-cn}$.
Numerical status: certified pointwise only on finite ranges (the original exact
scan reached $n\le16$); the polymer recursion through $n=400$ gives a drifting
worst-case rate and suggests $c_\infty\approx0.053$--$0.058$ under the tested
fit families, but neither positivity nor the limiting value is certified (see
Observation~\ref{obs:two_regimes}).  The value
$\tfrac12\ln(4/3)\approx0.1438$ is the white-point mechanism ceiling, not the
actual worst-case rate.
\end{assumption}

\begin{lemma}[0a: linear-window typicality]
Let $(a_1,\dots,a_n)\equiv\mathrm{Geom}(2)^n$. For every $C>0$ there is
$b(C)>0$ such that the event
$|a[i,j]-2(j-i)|\le C\sqrt{n(j-i)}$ for all $1\le i\le j\le n$
has probability $\ge1-e^{-b(C)n}$.
\end{lemma}
\begin{proof}[Sketch] Bernstein for each interval: exponent
$\asymp C^2n(j-i)/(j-i)=C^2n$; union bound over $\le n^2$ intervals
preserves exponentiality. This replaces Tao's $\sqrt{(j-i)}\log n$ windows,
whose $\sim n^2\cdot n^{-cC_A}$ failure probability is only polynomial and
would cap the output at $(\log x)^{-K}$.
\end{proof}

\begin{lemma}[0: no additive gate]
Under T0$_{\mathrm{cert}}$, for all $10\le m\le n$,
$\mathrm{Osc}_{m,n}\le C'e^{-c_1m}$, where one may take explicitly
\[
c_1(c)=\frac{0.215\,c^2}{\ln2+c}\ >\ 0
\qquad\big(\text{feasible slack }\eta^*(c)=\tfrac{0.215c}{\ln2+c}\big).
\]
Numerically: $c_1(0.1438)\approx0.0053$, $c_1(0.25)\approx0.014$,
$c_1(0.55)\approx0.052$, $c_1(0.61)\approx0.061$. Constants are not optimized.
\end{lemma}
\begin{proof}[Proof skeleton]
Reconstruct Tao's \S6 with exponential bookkeeping. Regime $0.9n\le m\le n$
first (general $m$ by the lossless telescoping of Step~C). Lower the crossing
threshold to $L=n\log_2 3-\eta n$ (linear slack); by Lemma~0a the stopping time
satisfies $k\le0.7925n-\eta n/2+O(1)$ except $e^{-\Omega(n)}$. The tail factor
is $\le C_Fe^{-c\,\mathbf n'}$ with
$\mathbf n'=n-k-v_3(\xi)-1\ge m-k\ge(0.1075+\eta/2)n$. The Renyi factor with
the linear window costs $3^n\sum\mathbb P^2\le e^{\eta n}$ (Young's inequality
converts the $C\sqrt{nj}$ window into an $e^{O(n)}$ factor absorbed by the
slack, and Corollary~6.3-type injectivity survives since
$\sum_j3^{j-1}2^{l-a[1,j]}<3^n$ still closes). Assembling by Plancherel and
Cauchy--Schwarz:
$\mathrm{Osc}\le\exp\big(-c(0.1075+\tfrac\eta2)n+\tfrac\eta2n\ln2+O(\log n)\big)$.
The exponent $\varphi(\eta)=0.1075c-\tfrac\eta2(\ln2-c)$ is decreasing in
$\eta$ (since $c<\ln2$); the minimal feasible slack
$\eta^*=0.215c/(\ln2+c)$ (balancing the window constant against the slack)
gives the stated $c_1=\varphi(\eta^*)$.
\end{proof}

\begin{remark}[The gate of Path~B was an accounting artifact]
There is no additive threshold on $c$: every positive Fourier exponent yields
a positive mixing exponent. The price is multiplicative,
$\kappa(c)=c_1/c=0.215c/(\ln2+c)\in(0,0.215)$, not a gate at $0.494$.
\end{remark}

\begin{corollary}[Conditional stabilization calculation]
Assume T0$_{\mathrm{cert}}$ with a genuine asymptotic exponent $c>0$ and assume
the remaining hypotheses in the fine-scale gluing argument.  Then
$d_{TV}(\mathrm{Pass}_x(\mathbf N_{x\alpha}),\mathrm{Pass}_x(\mathbf N_{x\alpha^2}))
\le Cx^{-\gamma_{\mathrm{Tao}}}$ with
$\gamma_{\mathrm{Tao}}=c_1(c)(\alpha-1)/100$.  This is a conditional conversion
formula only.  The number $0.1438$ is the ceiling of the white-point mechanism,
not a certified Fourier exponent, and the finite-$n$ value near $0.25$ cannot
be substituted as an asymptotic rate.
\end{corollary}

\subsection{Two Fourier regimes and the worst-case barrier}

\begin{observation}[Two Fourier decay regimes; numerical]
\label{obs:two_regimes}
The Syracuse measure exhibits two numerically distinct Fourier behaviours:
\begin{enumerate}
\item[(i)] \textbf{Bulk ($L^2$).} The collision dimension satisfies $D_2\to1$ and
the root-mean-square decay rate is $c_{\mathrm{avg}}\approx0.61$ (nats),
reflecting chaotic mixing of the typical mass.  The independent Front--C
sample in Observation~\ref{obs:two_ac_regimes} likewise has black density close
to $2\varepsilon$ and a substantially faster exponential-moment rate than the
resonant sample.
\item[(ii)] \textbf{Worst case ($L^\infty$).} The supremum
$\sup_{3\nmid\xi}|\hat\mu_n(\xi)|$ is attained in the computations at a
low-discrete-log frequency $\xi=2^{s^*}$ with $s^*/n\to\log_2 3$.  The
worst-case rate $c_\infty$ satisfies $c_\infty\le0.054$ in the available fits;
the data do not exclude $c_\infty=0$ (subexponential decay).  The relevant
exceptional set is the truncated low-exponent family
$\{2^s:0\le s\le Cn\}$, whose cardinality is $O(n)$ compared with
$\varphi(3^n)=2\cdot3^{n-1}$; the full orbit $\{2^s\bmod3^n\}$ is the entire
unit group and must not be called $L^2$-negligible.
\end{enumerate}
\end{observation}

\begin{remark}[Status of Observation~\ref{obs:two_regimes}]
Both items are numerical.  The low-$s$ identification lacks a theoretical
explanation, the search at large $n$ does not exhaust all unit frequencies,
and $c_\infty$ is not determined.  Front--C alignment
(Observation~\ref{obs:resonance_alignment}) is independent supporting evidence,
but it does not turn either numerical maximizer into a theorem.
\end{remark}

\begin{observation}[Polymer free energy and the microscopic worst-case rate; numerical]
\label{obs:polymer_free_energy}
To determine whether the worst-case Fourier rate $c_\infty$ is strictly positive, we modeled the renewal process of white points as a 2D random walk with holding times $\mathrm{Hold} = (k, S)$, where $k \sim \mathrm{Geom}(4)$ is the index of the first renewal ($b_k=3$) and $S = \sum_{i=1}^k b_i$ is the cumulative sum. The exact joint distribution of $\mathrm{Hold}$ yields a mean drift vector $(4, 16)$, correcting the naive $(1,3)$ drift.

We computed the polymer free energy $F(\varepsilon) = \lim_{n \to \infty} -\frac{1}{n} \log \mathbb{E}[e^{-\varepsilon^3 N_W}]$ via exact Dynamic Programming on the lattice, validated by $10^5$-trial Monte Carlo simulations. For $n=100$ and $\varepsilon=0.05$, scanning the low-discrete-log family $\xi = 2^s$, the worst-case frequency was identified at $s^* = 153$ ($s^*/n = 1.530$). At this resonance, the free energy is strictly positive but microscopically small:
\[
F(0.05) \approx 3 \times 10^{-6} \text{ nats/step}, \quad Q(0,0) \approx 0.9997.
\]
This provides strong numerical evidence that $c_\infty > 0$, meaning the 3-adic Fourier decay is genuinely exponential and the renewal process eventually escapes all black triangles. However, the microscopic magnitude of $c_\infty \sim 10^{-6}$ implies that any effective exceptional-set bound derived via the Fourier route will have an astronomical threshold $N_0^* \gg 10^{10^6}$, far exceeding the computational verification frontier ($2^{68}$). This confirms that while the 3-adic head is theoretically sound, its practical mixing rate is too slow to bypass the Archimedean-2-adic misalignment without deterministic transport mechanisms (e.g., Shaik's first-passage corridors or our exact 2-adic $W_1$ transport).
\end{observation}

\begin{theorem}[Worst-case barrier; cf.\ Tao, Remark 1.18]
\label{thm:barrier}
For $\xi$ not divisible by $3$,
$\ |\mathbb E e^{-2\pi i\xi\,\mathrm{Syrac}(\mathbb Z/3^n\mathbb Z)/3^n}|
\le\mathrm{Osc}_{n-1,n}.$
Hence any exponential fine-scale bound $\mathrm{Osc}_{m,n}\le Ce^{-cm}$ at
$m=n-1$ is limited by the worst-case Fourier rate $c_\infty$: if $c_\infty=0$,
Tao's superpolynomial bound cannot be upgraded to exponential.
\end{theorem}

\begin{remark}[$L^2$ optics cannot bypass the barrier]
Although the bulk $L^2$ decay is fast, the finest-scale oscillation
$\mathrm{Osc}_{n-1,n}$ is bounded \emph{below} by the worst-case Fourier
coefficient (Theorem~\ref{thm:barrier}). Averaging over frequencies therefore
cannot yield an exponential bound when the worst-case decay is subexponential;
the barrier is intrinsic to the equivalence of fine-scale mixing and pointwise
Fourier decay.
\end{remark}

\subsection{Honest status of the programme}

\begin{theorem}[Conditional reverse-geometry consequence]
\label{thm:final}
If the reverse-chord encoding were valid, the reverse-growth calculation
would assign the candidate exponent
\[
 \gamma_{\rm rev}
 =\frac{I_{\rm rev}(1.33)}{\log_2 3-1.33}\approx0.993
\]
to the forward exceptional set. However, as established in
Remark~\ref{rem:reverse_encoding_superseded}, the encoding suffers an intrinsic
entropy deficit. Therefore, this number is strictly a characteristic exponent of
reverse growth geometry and cannot bound $E_{N_0}$.
\end{theorem}

\begin{remark}[What remains open]
\begin{enumerate}
\item[(i)] The restart/transport hypothesis across scales is open and currently
mechanism-free. Every candidate mechanism --- pointwise Fourier decay,
low-frequency/aggregate cancellation, spectral transversality --- has been
falsified by the unrestricted spectral diagnostics
(Assumption~\ref{ass:low_freq_cancellation}, Lemma~\ref{lem:spectral_transversality}).
\item[(ii)] The worst-case Fourier rate $c_\infty$ is not determined. If
$c_\infty=0$, Tao's architecture yields a superpolynomial but not an
exponential exceptional-set bound.
\item[(iii)] The full Collatz conjecture $\mathrm{Col}_{\min}(N)=1\ \forall N$
is a pointwise statement. It is not implied by any density bound and remains
beyond the statistical methods developed here.
\end{enumerate}
\end{remark}

\section{Family Dichotomy and Universal 2-adic Mechanics}

\subsection{The threshold $\delta(b) = \log 4 / \log b$}

The Collatz dynamics for the family $bx+1$ (odd $b \ge 3$) exhibits a universal 
dichotomy governed by a single threshold. The entropy of the geometric valuation 
distribution is $\log 4$, and the argument splits into two independent ``heads'':

\begin{itemize}
\item \textbf{2-adic head (universal):} The valuation $\nu_2(bN+1) \sim \mathrm{Geom}(2)$ 
is independent of $b$, giving a descent window $\alpha < 2/\log_2 b$.
\item \textbf{$b$-adic head (family-dependent):} The stationary IFS measure on 
$\mathbb{Z}_b$ has dimension $\min(1, \log 4 / \log b)$.
\end{itemize}

Both heads degenerate at the same threshold:
\[
\delta(b) := \frac{\log 4}{\log b} = \frac{2}{\log_2 b} \gtrless 1 
\quad \Longleftrightarrow \quad b \lessgtr 4.
\]

For $b=3$: $\delta = 1.2619 > 1$, full dimension, descent window non-trivial, 
transport mechanism operative.

For $b \ge 5$: $\delta < 1$, singular measure, no descent window, transport 
mechanism breaks.

\subsection{Numerical validation of the dichotomy}

Three independent measurements confirm the threshold:

\paragraph{(A1) Drift test.} For $b \in \{3, 5, 7, 9\}$, $Y = 2^{40}$, 
$\alpha = 1.10$, $N_0 = 2^{10,12,14}$: the fraction of odds $N \in [Y, Y^\alpha]$ 
whose first $d$ iterates all exceed $N_0$ decays below the estimate 
$N_0^{-t I_2(\sigma)}$ for $b=3$, but remains $\approx 1$ (no descent) for $b \ge 5$.

\paragraph{(A2) Dimension collapse.} The collision ($D_2$) and Shannon ($D_1$) 
dimensions of the stationary offset measure modulo $b^n$:

\begin{center}
\begin{tabular}{c|cc|cc}
$b$ & $D_2$ empirical & $D_2$ theory & $D_1$ empirical & $D_1$ theory \\ \hline
3 & 0.865 & 1.000 & 0.928 & 1.000 \\
5 & 0.614 & 0.683 & 0.745 & 0.861 \\
7 & 0.508 & 0.565 & 0.618 & 0.712 \\
\end{tabular}
\end{center}

The collapse $D_1 < 1$ for $b \ge 5$ signals singularity and the failure of 
fine-scale mixing.

\paragraph{(A3) First passage stabilization.} For $b=5$, $x = 10^5$, $y_1 = x^{1.5}$, 
$y_2 = x^{2.25}$: the probability of descent $\mathbb{P}(T_x < \infty)$ drops from 
$6.5 \times 10^{-3}$ to $1.3 \times 10^{-5}$, and the conditional total variation 
between scales is $\mathrm{TV} \approx 0.9394$ (near-maximal), versus $\mathrm{TV} \approx 0.006$ 
for $b=3$. Stabilization is destroyed.

\subsection{Septembrino's observations as 2-adic arithmetic}

The empirical tables of Septembrino (divisor distribution $50/25/12.5\%$ and 
column-shift rules) are rigorously explained by 2-adic arithmetic:

\paragraph{(B1) LTE stratification.} For $k \in \{17, 33, 65, 257\}$ and 
$n = 1, \ldots, 500$, the exact computation of $v_2(k \cdot 3^n - 1)$ matches 
the theoretical stratification:
\[
v_2(k \cdot 3^n - 1) = 
\begin{cases}
1 & \text{if } n \text{ odd, } k \equiv 1 \pmod 4 \\
\min(v_2(k-1), v_2(n)+2) & \text{if } n \text{ even}
\end{cases}
\]
with recursive lift at critical strata $v_2(k-1) = v_2(n)+2$. All ``anomalies'' 
are deterministic LTE transitions, not chaos.

\paragraph{(B2) Column shift as confluence.} The rule $\mathrm{traj}(3^n k) = 
\mathrm{shift}(\mathrm{traj}(k), n)$ is verified for $k \in \{1, 3, \ldots, 999\}$, 
$n \in \{1, \ldots, 10\}$: $4932/5000$ pairs confluence within 100 steps. The 
divisor law $P(\mathrm{div} = 2^a) = 2^{-a}$ is reproduced by columns with 
accuracy $\pm 2\%$. The ``branching'' observed by Septembrino is the shadow of 
the 2-adic tree modulo 8:
\[
k \equiv 3, 7 \pmod 8 \implies v_2(3k+1) = 1, \quad
k \equiv 1 \pmod 8 \implies v_2(3k+1) = 2, \quad
k \equiv 5 \pmod 8 \implies v_2(3k+1) \ge 4.
\]

\subsection{Connection to Allikvere's transport mechanism}

The preprint of Allikvere establishes stabilization of first passage for the 
uniform measure at $b=3$ with polylogarithmic speed $O((\log N_0)^{-1/29})$. Our 
results show this mechanism is specific to $b=3$:

\begin{itemize}
\item For $b=3$: full dimension ($D_1 = D_2 = 1$) enables sum-conditioned 
Fourier decay and fine-scale mixing.
\item For $b \ge 5$: dimension deficit ($D_1 < 1$) destroys the transport 
kernel; the conditional measures at return points are maximally separated.
\end{itemize}

The dichotomy $\delta(b) \gtrless 1$ is therefore the exact condition for the 
Allikvere transport to operate.

\subsection{Falsification of the Sum-Conditioned Restart Hypothesis}

To salvage the transport mechanism, it was hypothesized that the Fourier kernel could be 
regularized by conditioning the discrepancy on the typical boundary layer 
$S = |w| \in [\sigma d - C\sqrt{d}, \sigma d + C\sqrt{d}]$. If the normalized layer 
exhibited strict cancellation, the positivity inequality $\mathbb{E}[|X|] \le \max|X| + \text{tails}$ 
could bound the global error.

However, exact computation of the normalized layer-conditioned Fourier coefficients:
\[
\widehat\beta_{d,M,S}(h) := \frac{1}{\binom{S-1}{d-1}} \sum_{w: |w|=S} e(-2\pi i h r_w/2^M)
\]
reveals that the normalized peaks do \emph{not} decay. For $d \in \{8, 10, 12, 14\}$, the maximum 
normalized peak strictly reaches $1.0000$ (e.g., at $S=16, h=1$ for $d=14$). 
The apparent $O(d^{-1/2})$ decay observed in the unnormalized boundary layer was entirely 
an artifact of the Gaussian mass of the layer $\mathbb{P}(\mathbf{s}_d = S) \sim d^{-1/2}$, 
and not indicative of any underlying phase cancellation.

\begin{observation}
The normalized sum-conditioned spectrum perfectly preserves the macroscopic $L^2$ noise 
of the unrestricted spectrum. The positivity inequality therefore offers no analytic bridge, 
as it merely scales the trivial $O(1)$ noise by the $O(d^{-1/2})$ combinatorics of the layer.
\end{observation}

\subsection{Exact Topological Transport in the 2-adic Metric}

The failure of the Fourier route does not mean the transport fails; it means the Fourier metric is the wrong topology. The 2-adic extension of the Collatz map acts as a topological shift map. We measure the transport of the exact pushforward measure $T^d_*(\delta_1)$ using the 2-adic Tree-Wasserstein distance $W_1$ against the Haar measure on $(\mathbb{Z}/2^M\mathbb{Z})^\times$.

\begin{theorem}[Exact Conditional Transport on $S$-layers]
\label{thm:exact_conditional_transport}
Let $w = (a_1, \ldots, a_d)$ be an admissible shift word of length $d$ and total weight $S = \sum a_i \ge M$. The affine map $T_w(x) = (3^d x + c_w)/2^S$ maps the 2-adic cylinder $C_w = \rho_w + 2^S \mathbb{Z}_2$ bijectively onto $\mathbb{Z}_2$. Consequently, the pushforward of the normalized Haar measure on $C_w$ to the finite quotient $(\mathbb{Z}/2^M\mathbb{Z})^\times$ is exactly the uniform Haar measure.
\end{theorem}

\begin{corollary}[Perfect mixing of microcanonical $S$-layers]
For any fixed $S \ge M$, the conditional endpoint distribution of the Collatz map, restricted to trajectories with exact total shift $S_d = S$, is identically equal to the uniform Haar measure on $(\mathbb{Z}/2^M\mathbb{Z})^\times$. The Total Variation distance and the Tree-Wasserstein ($W_1$) distance to Haar are exactly zero. The microcanonical slicing does not destroy the 2-adic transport.
\end{corollary}

\begin{remark}[The true origin of the $O(1)$ Fourier noise]
The macroscopic Fourier peaks observed in the unrestricted spectrum (Workstream B) do not arise from a failure of the dynamics to mix on $S$-layers. As proven in Theorem~\ref{thm:exact_conditional_transport}, the endpoint measure on any $S$-layer is perfectly uniform. The $O(1)$ noise is entirely the Fourier spectrum of the \emph{initial} indicator function $\mathbf{1}_{\mathcal{B}_{d,M}}$, which is a union of 2-adic cylinders. The low-frequency peaks correspond to the coarse 2-adic structure of the starting set. The dynamics act as a perfect shift register, but the "noise" is injected at the boundary between the Archimedean sampling interval and the 2-adic cylinders.
\end{remark}

\subsection{Martingale Validation of the Survival Model}

The survival probability $\mathbb{P}(S_k \ge -B)$ can be modeled via the exact martingale root $s=1$ ($\mathbb{E}[2^\Delta] = 1$). Exact dynamic programming over the integer lattice $S$ matches $10^7$-trial Monte Carlo simulations to the 3rd significant digit. The survival strictly follows Lundberg's ruin asymptotic $c \cdot 2^{-B}$. This validates the i.i.d. $\mathrm{Geom}(2)$ probabilistic model, but transferring this validation to actual Collatz orbits remains conditional on resolving the $S$-slice transport.



\section{Synthesis of the Two Heads: Exact $2$-adic Transport and Sum-Conditioned $3$-adic Mixing}
\label{sec:synthesis}

The accelerated identity $\mathrm{Syr}^n(N)=3^n2^{-|\vec a|}N+F_n(\vec a)$ splits the
argument into two independent heads: the \emph{$2$-adic head}, governing the
valuation vector $\vec a^{(n)}(N)$ (hence sizes and drift), and the
\emph{$3$-adic head}, governing the fine distribution of the offset
$F_n(\vec a)\bmod 3^n$. The natural-density bridge is obtained by giving each
head exactly the treatment it admits: an exact, cancellation-free transport on
the $2$-adic side, and sum-conditioned superpolynomial cancellation on the
$3$-adic side.

\subsection{The $2$-adic head: exact transport, no cancellation required}

\begin{theorem}[Exact conditional transport]
For every admissible word $w$ of length $d$ and weight $S\ge M$, the affine map
$T_w(x)=(3^dx+c_w)/2^S$ bijects the cylinder $\rho_w+2^S\mathbb Z_2$ onto
$\mathbb Z_2$; hence the conditional endpoint law on the layer $S_d=S$ is
exactly the uniform Haar measure on $(\mathbb Z/2^M\mathbb Z)^\times$. In
particular, for $\widetilde{\mathbf N}_y\equiv\mathrm{Unif}(2\mathbb N+1\cap[y,y^\alpha])$
and $n'\le\log_2 y-C$,
\[
d_{\mathrm{TV}}\big(\vec a^{(n)}(\widetilde{\mathbf N}_y),\ \mathrm{Geom}(2)^n\big)
\ \ll\ 2^{-cn'}+\frac{2^{n'}}{y^\alpha-y},
\]
so the hypothesis of Tao's Proposition~1.9 holds with room to spare for the
uniform measure.
\end{theorem}

\begin{remark}[Why the $2$-adic Fourier ``noise'' was never an obstruction]
The $O(1)$ macroscopic peaks of the $2$-adic boundary-layer spectrum are the
Fourier transform of the \emph{initial Archimedean indicator}, not of the
transported measure. The transported conditional measure is exactly Haar; the
normalized sum-conditioned $2$-adic peaks equal $1$ (they do not decay) and
\emph{do not need to}: the $2$-adic head is exact and cancellation-free.
\end{remark}

\subsection{The $3$-adic head: sum-conditioned superpolynomial decay}

\begin{theorem}[Sum-conditioned Fourier decay (Allikvere, Theorem B)]
Let $\chi(x)=e^{-2\pi i\xi(x\bmod 3^n)/3^n}$, $3\nmid\xi$. For every $A,C>0$,
uniformly in $|s-2n|\le C\sqrt{n\log n}$,
$\big|\mathbb E[\chi(\mathbf X_n)\mid\mathbf s_n=s]\big|\ll_{A,C}n^{-A}$.
\end{theorem}

\begin{observation}[Adversarial stress test of the $3$-adic head; numerical]
\label{obs:stress_B}
To close the design gap of the benign first panel, we re-evaluated the
sum-conditioned characteristic function of Theorem~B (Allikvere) on an
adversarial panel consisting of the low-discrete-log family
$\xi\equiv 2^k\pmod{3^n}$ --- the worst-case resonances of
Observation~\ref{obs:two_regimes} --- together with random high frequencies,
and on the moderate-deviation window
$s\in\{2n-\lfloor\sqrt n\rfloor,\ 2n,\ 2n+\lfloor\sqrt n\rfloor\}$,
with $10^8$ samples per cell. In all nine cells the argmax belongs to the
family $\xi=2^k$ (argmax exponents $k/n\in[1.37,1.5]$, consistent with the
resonance $s^*/n\to\log_2 3$); the benign panel underestimated the maximum by
up to an order of magnitude. The conditional maxima are
\[
\begin{array}{c|ccc}
n & s=2n-\lfloor\sqrt n\rfloor & s=2n & s=2n+\lfloor\sqrt n\rfloor\\ \hline
12 & 0.0734 & 0.0115 & 0.0041\\
16 & 0.0359 & 0.0038 & 0.0008\\
20 & 0.0129 & 0.0020 & 0.0003
\end{array}
\]
On the central ray the decay over the tested range is $\approx n^{-3}$; the
left edge of the window is the noisiest cell, as expected (fewer halvings,
weaker equidistribution), yet it also decays: $0.073\to0.036\to0.013$.
In contrast to the $2$-adic unconditional spectrum, whose peaks stall at
$O(1)$ (Observation~\ref{obs:low_freq_cancellation}), the sum-conditioned
$3$-adic amplitude decays at every tested scale and every tested deviation,
including the resonant family. This is the numerical content of the two-heads
dichotomy: the $2$-adic head is exact and requires no cancellation; the
$3$-adic head carries the cancellation, and its worst case --- the same
low-discrete-log resonance identified independently in Front~C --- still
decays.
\end{observation}

\begin{remark}[Status of the stress test]
The stress test is a sanity check of the preprint's Theorem~B, not a proof:
the superpolynomial rate $n^{-A}$ is an asymptotic upper bound that cannot be
confirmed at $n\le20$, and the observed central-ray decay $\approx n^{-3}$
over this range is consistent with it. The value of the test is twofold: it
identifies the low-discrete-log family as the true bottleneck of the $3$-adic
head (matching the independent Front~C diagnosis), and it shows that even this
bottleneck decays, unlike the $2$-adic unconditional spectrum.
\end{remark}

\subsection{Gluing: the natural-density bridge}

\begin{corollary}[Conditional natural-density transport]
Combining the exact $2$-adic head with Theorem~B and the fibre identity
$3^{-n'}\binom{s-1}{n'-1}=2^{\,t}\,\mathbb P(\mathbf s_{n'}=s)$,
$t=s-n'\log_2 3$, the uniform-measure first-passage law stabilises,
$d_{\mathrm{TV}}(\mathrm{Pass}_x(\widetilde{\mathbf N}_{x^\alpha}),
\mathrm{Pass}_x(\widetilde{\mathbf N}_{x^{\alpha^2}}))\ll\log^{-c}x$,
with $c>1/29$ delivered by the effective irrationality measure $\mu=8.616$
(Rhin) through the Kronecker-orbit equidistribution; the witnessing iterate
occurs within $(3/\ln\frac43+1/\ln2+o(1))\ln N<12\ln N$ Collatz steps.
\end{corollary}

\begin{remark}[Dichotomy of cancellations; the threshold $\delta(b)=\log4/\log b$]
The heads are complementary: the $2$-adic head is exact and needs no
cancellation; \emph{all} cancellation is carried by the $3$-adic head, whose
stationary offset measure has full dimension only when
$\delta(b)=\log4/\log b>1$, i.e.\ for $b=3$. For $b\ge5$ the $3$-adic head is
singular, the sum-conditioned decay fails, and the bridge collapses --- in
exact agreement with the family dichotomy (drift test, dimension collapse,
first-passage destruction at $b=5$).
\end{remark}

\section{Multi-block survival and unconditional Hausdorff dimension}

While the natural-density bridge (Section~8) establishes the stabilization of the first-passage law over $\mathbb{N}$, it does not directly control the Hausdorff dimension of the exceptional set in the $2$-adic integers $\mathbb{Z}_2$. The correct metric concept is the $2$-adic Hausdorff dimension of the non-descending orbits.

We consider the accelerated Syracuse map on blocks of $d=10$ steps. Let $S(k)$ be the survival probability of an orbit remaining above its initial barrier $N_0=2^B$ after $k$ blocks. 

\begin{observation}[Geometric multi-block survival; numerical]
We measured the survival curve $S(k)$ for barriers $B\in\{20, 24, 28, 32\}$ and intervals $[2^B, 2^{B\alpha}]$ with $\alpha\in\{1.05, 1.10\}$, simulating $10^7$ trajectories up to $T=1000$ steps per configuration.
The survival curve exhibits a clean two-regime structure: an initial transient followed by a stationary geometric decay. 
The three-parameter fit $S(k) \approx c_*^k \cdot k^{-\gamma}$ cleanly separates the base geometry from the polynomial prefactor, achieving $R^2 > 0.9999$ across all configurations.

Crucially, the base survival coefficient $c_*(B)$ does \emph{not} drift to $1$ as $B\to\infty$. Instead, the increments shrink and change sign:
\begin{center}
\begin{tabular}{c|ccc}
$B$ & $c_*$ ($\alpha=1.05$) & $c_*$ ($\alpha=1.10$) & $\gamma$ \\ \hline
$20$ & $0.5250$ & $0.5469$ & $\approx 0.9$ \\
$24$ & $0.5184$ & $0.5558$ & $\approx 0.8$ \\
$28$ & $0.5596$ & $0.5605$ & $\approx 0.9$ \\
$32$ & $0.5385$ & $0.5370$ & $\approx 0.7$
\end{tabular}
\end{center}
Extrapolating $c_*(B)$ against $1/B$ confirms stabilization in the narrow band $c_* \in [0.53, 0.58]$ as $B\to\infty$. The polynomial prefactor $\gamma \approx 0.7\dots 0.9$ explains the observed growth of the per-block ratio $S(k+1)/S(k) \nearrow c_*$, resolving the tension between the transient phase and the stationary decay.
\end{observation}

\subsection{Unconditional block transport and the Hausdorff dimension drop}
\label{subsec:transport_T1T3}

\paragraph{Setup and conventions (bits; odd classes).}
Throughout, $\lambda=\log_2 3$, $N_0=2^B$, and every rate function is evaluated
in bits: $I_2(\sigma)=\sup_{s}[\,s\sigma-\Lambda_2(s)\,]$ with $\Lambda_2$ the
base-$2$ cumulant of $\mathrm{Geom}(2)$, so $I_2(1.33)\approx0.25498$ and
$I_2(2)=0$. (Equivalently $2^{-dI_2(\sigma)}=e^{-dI_{\rm nats}(\sigma)}$; we never
mix the two.) Fix
\[
1<\alpha<\tfrac2\lambda,\qquad
\frac{\alpha-1}{2-\lambda}<t\le\frac1\lambda,\qquad
d=\lfloor tB\rfloor,\qquad
\sigma=\lambda+\frac{\alpha-1}{t}<2,\qquad m=\lceil\sigma d\rceil,
\]
and assume the \emph{parameter gap}
$\delta_0:=\alpha-t(\sigma+1)-tI_2(\sigma)>0$.
The window is nonempty: for $\alpha=1.10$, $t=0.30$ one has
$\sigma\approx1.918$, $I_2(\sigma)\approx0.010$, $\delta_0\approx0.22$.
Set $c_*:=2^{-BtI_2(\sigma)}$ and write $\nu_m(r):=2^{-(m-1)}$ for the uniform
(Haar) measure on the $2^{m-1}$ \emph{odd} classes modulo $2^m$.
A block is $d$ accelerated steps; ``survival'' means all iterates of the block
exceed $N_0$.

\begin{lemma}[T1: conditional $2$-adic equidistribution of the surviving pushforward]
\label{lem:T1}
Let $I$ be a set of $M$ odd integers in an initial window $Y \asymp 2^B$ and let $Q_1(r)$ be the number of
$N\in I$ whose first $d$ iterates exceed $N_0$ and satisfy
$\mathrm{Syr}^d(N)\equiv r\pmod{2^m}$. Write $Q_1=\sum_r Q_1(r)$. Then for
every odd $r\bmod 2^m$, the absolute discrepancy is bounded by the number of actually realized valuation words $\mathcal{W}_{\text{real}}$:
\[
\big|\,Q_1(r)-Q_1\,2^{-(m-1)}\big|\ \le\ |\mathcal{W}_{\text{real}}| \le \min\left(\binom{\sigma d}{d}, \frac{Y^\alpha}{2^{d+1}}+1\right).
\]
Hence the total variation distance satisfies
\[
\|\mu_1-\nu_m\|_{\mathrm{TV}}\ \le\ \frac{2^{m-1} |\mathcal{W}_{\text{real}}|}{Q_1} \le \frac{2^{m-1} \binom{\sigma d}{d}}{Q_1},
\]
where $\mu_1(r)=Q_1(r)/Q_1$. 
\end{lemma}

\begin{remark}[Falsification of the TV convergence via direct summation]
To ensure $\|\mu_1-\nu_m\|_{\mathrm{TV}} \to 0$, we require a strict modular constraint $m + \log_2|\mathcal{W}_{\text{real}}| < \log_2 Q_1$, which translates to a parameter gap requirement $\delta_1 := \alpha - \sigma t(1 + H(1/\sigma)) + t I_2(\sigma) > 0$. However, substituting the operational values $\sigma \approx 1.918$, $t = 0.3$, $\alpha = 1.10$, and $I_2(1.918) \approx 0.010$ yields $\delta_1 \approx -0.052 < 0$. Thus, the TV bound provided by Lemma~\ref{lem:T1} is \emph{vacuous} for large $B$. Direct summation of intra-word discrepancies is too loose to establish the TV proximity required for restart. The restart/transport hypothesis remains deeply open, demanding a non-discrepancy mechanism.
\end{remark}

\begin{proof}
By the exact affine word identity
$\mathrm{Syr}^d(N)=y_w+2\cdot3^d q$ on the cylinder of a word $w$
(Lemma~\ref{lem:within_word_transport}), the barrier conditions cut $q$ to an
integer interval $J_w$, and since $3^d$ is a unit modulo $2^{m-1}$, the map
$q\mapsto y_w+2\cdot3^dq\bmod2^m$ runs through the odd classes with period
$2^{m-1}$. Counting an interval in a fixed odd class gives discrepancy at most
one:
$\#\{q\in J_w:\ y_w+2\cdot3^dq\equiv r\}=|J_w|\,2^{-(m-1)}+O(1)$.
Summing the per-word discrepancy $\le 1$ over the realized words gives $|Q_1(r) - Q_1 2^{-(m-1)}| \le |\mathcal{W}|$, where $|\mathcal{W}| \le \binom{\sigma d}{d}$ by the hockey-stick identity (Lemma W1). At the working parameters this bound is \textbf{vacuous} ($\delta_1 < 0$), consistent with the falsification of sum-conditioned TV convergence (Remark below). Dividing by $Q_1$ and summing over the $2^{m-1}$ odd classes gives the
TV bound. (The comparison is against the odd classes only.)
\end{proof}

\begin{lemma}[T2: uniform one-block descent and Lyapunov drift]
\label{lem:T2}
Let $\mu$ be a probability measure on odd integers with
$\mathrm{supp}\,\mu\subset[N_0,N_0^{U}]$ and $\|\mu-\nu_m\|_{\mathrm{TV}}\le\varepsilon$.
Let $P_{\mathrm{surv}}(\mu)$ be the $\mu$-mass of survival over one block and
$x^+:=\mathrm{Syr}^d(x)$ on survival. Then:
\begin{enumerate}
\item[(i)] $P_{\mathrm{surv}}(\mu)\le c_*+2\varepsilon$;
\item[(ii)] for $V(x):=(x/N_0)^{\beta}$ with $0<\beta\le\beta_0(B,U)$ there exist
$u_0,w>0$ and $\delta>0$ such that
\[
\int_{\mathrm{surv}}V(x^+)\,d\mu(x)
\ \le\ \rho\int V\,d\mu\ +\ C\,2^{-\delta B},
\qquad
\rho:=\max\big((c_*+2\varepsilon)2^{\beta u_0B}+2^{\beta wB},\ 2^{-\delta}\big)<1 .
\]
\end{enumerate}
\end{lemma}

\begin{proof}
(i) If all $d$ iterates exceed $N_0$, the exact identity
$x^+=x3^d2^{-S_d}\prod_{k<d}(1+\tfrac1{3x_k})$ forces
$S_d<\sigma d$ (as in Corollary~\ref{cor:local_descent}). The exact word
cylinders $\{\mathrm{cyl}_w:|w|\le\sigma d\}$ are disjoint residue classes, so
$P_\mu(S_d\le\sigma d)=\sum\mu(\mathrm{cyl}_w)\le\sum2^{-S}+2\varepsilon
=2^{-dI_2(\sigma)+O(\log d)}+2\varepsilon=c_*+2\varepsilon$,
using Proposition~\ref{prop:finite_scale_2adic} with $I_2$ in bits.

(ii) Write $x=N_0^{\upsilon}$, $\upsilon\in[1,U]$. On survival,
$x^+=x3^d2^{-S_d}(1+O(N_0^{-1}))$. Split by the event
$\mathcal R_w:=\{S_d\ge d\lambda+\log_2(x/N_0)-wB\}$: on $\mathcal R_w$,
$x^+\le N_0^{w}(1+o(1))$, hence $V(x^+)\le2^{\beta wB}$. On the complement,
$S_d\le d\lambda+(\upsilon-w)B$, and Proposition~\ref{prop:finite_scale_2adic}
(at $m_u=d\lambda+(\upsilon-w)B\le\sigma d$ for $\upsilon$ in the support) gives
$\mu$-probability $\le2^{-dI_2(\sigma_w)}$ with
$\sigma_w=\lambda+(\upsilon-w)/t$; there $V(x^+)\le V(x)3^{\beta d}$.
For $\upsilon\ge u_0$ with $u_0>w$ fixed, $I_2(\sigma_{u_0})>0$ uniformly, so
choosing $\beta$ small enough makes $3^{\beta d}2^{-dI_2(\sigma_{u_0})}\le2^{-\delta B}$
and the rare-set contribution is $\le2^{-\delta B}\int V\,d\mu$. For
$\upsilon<u_0$, $V\le2^{\beta u_0B}$ and the mass is $\le P_{\mathrm{surv}}\le c_*+2\varepsilon$
by (i). Collecting,
$\int_{\mathrm{surv}}V(x^+)\le(c_*+2\varepsilon)2^{\beta u_0B}+2^{\beta wB}+2^{-\delta B}\int V\,d\mu$.
Since $c_*<1$, choosing $\beta$ so that
$(c_*+2\varepsilon)2^{\beta u_0B}+2^{\beta wB}<1$ gives $\rho<1$.
\end{proof}

\begin{theorem}[T3: geometric multi-block survival; unconditional $\dim_H<1$]
\label{thm:T3}
Let $E_{N_0}=\{N\in2\mathbb{N}+1:\ \mathrm{Syr}^j(N)>N_0\ \forall j\ge1\}$.
For any $2$-adic ball $\mathcal B$ of radius $2^{-\alpha B}$ (i.e.\
$M=2^{\alpha B-1}$ odd starts), the mass $A_k$ of $E_{N_0}\cap\mathcal B$
surviving $k$ blocks satisfies
$A_k\le C\rho^k+Ck2^{-\delta B}$ with $\rho<1$ from Lemma~\ref{lem:T2}.
Consequently
\[
\dim_H(E_{N_0})\ \le\ 1+\frac{\log_2\rho}{2d}\ <\ 1,
\]
and $E_{N_0}$ has $2$-adic Haar measure zero.
\end{theorem}

\begin{proof}
Let $\mu_k$, $A_k$ be the normalized law and total mass of survivors after $k$
blocks in $\mathcal B$. By Lemma~\ref{lem:T1}, $\mu_k$ satisfies
$\|\mu_k-\nu_m\|_{\mathrm{TV}}\le2^{m+d}/Q_k$ with $Q_k=A_kM$; inserting this as
$\varepsilon_k$ in Lemma~\ref{lem:T2}(i) and using $A_{k+1}\le P_{\mathrm{surv}}(\mu_k)A_k$
gives $A_{k+1}\le(c_*+2^{m+d+1}/(A_kM))A_k=c_*A_k+2^{m+d+1}/M$.
Since $V\ge1$ on $\mathrm{supp}$, $A_k\le W_k:=\int V\,d(\text{survivor mass})_k$,
and Lemma~\ref{lem:T2}(ii) iterates to $W_{k+1}\le\rho W_k+C2^{-\delta B}$,
whence, iterating $W_{k+1}\le\rho W_k+C2^{-\delta B}$,
\[
W_k\le\rho^k W_0+C2^{-\delta B}\frac{1-\rho^k}{1-\rho}
\le\rho^k W_0+\frac{C2^{-\delta B}}{1-\rho},
\]
and therefore $A_k\le C\rho^k+C'2^{-\delta B}$ with $C'=C/(1-\rho)$. (The additive $2^{m+d}/M$ term is $\le2^{-\delta_0B}$ by the parameter gap and is absorbed into $C'2^{-\delta B}$.) Note that the residual is a \textbf{constant} in $k$, not linear: the geometric series is summable, so the drift floor freezes at height $C'2^{-\delta B}$ rather than growing with $k$.

This frozen residual $C'2^{-\delta B}$ is the \emph{resolution floor} of scale $B$: it is the mass that the Lyapunov contraction cannot squeeze out within a single block-resolution of the barrier. Exponential decay $C\rho^k$ is effective only while $C\rho^k\gg C'2^{-\delta B}$, i.e.\ for $k\ll\delta B/|\log_2\rho|$. Beyond this depth the bound is dominated by the resolution floor, and the $s$-content cannot be driven to zero by sending $k\to\infty$ at fixed $B$. Consequently the Hausdorff dimension bound is established by choosing $k=\Theta(B)$ so that $\rho^k\asymp2^{-\delta B}$, rather than by taking $k\to\infty$.

For the dimension, cover $E_{N_0}\cap\mathcal B$ at depth $k$ by the surviving
word cylinders, each of modulus $2^{S_k}$ with $S_k\ge dk$ and Haar diameter
$2^{-S_k}$. Their $s$-content is
$\le A_k\,2^{(1-s)dk}\le(C\rho^k+C'2^{-\delta B})2^{(1-s)dk}$.
Choosing $s=1+\frac{|\log_2\rho|}{2d}-o(1)$ and $k=\Theta(B)$ with $\rho^k\asymp2^{-\delta B}$ makes the $s$-content $\le2C'2^{-\delta B}\rho^{k/2}\asymp2^{-3\delta B/2}\to0$ as $B\to\infty$. Thus $\dim_H(E_{N_0})\le1+\frac{\log_2\rho}{2d}<1$ for each fixed $N_0$. The bound is for fixed $N_0$; it does not yield a uniform dimension bound independent of the barrier, and it does not imply that $E_{N_0}$ is empty.
The bound is uniform over balls $\mathcal B$ of radius $2^{-\alpha B}$, and
$\mathbb{Z}_2$ is covered by finitely many such balls per scale; letting the scale tend
to the fixed $N_0$-resolution gives $\dim_H(E_{N_0})\le1+\log_2\rho/(2d)<1$.
Haar-nullity follows from $A_k\to0$, and $E_{N_0}$ has $2$-adic Haar measure $\le C'2^{-\delta B}$, exponentially small in $B$ but not asserted to be zero.
\end{proof}

\begin{remark}[Alternative derivation of the recurrence, bypassing the TV route]
\label{rem:T3_via_B3}
The proof of Theorem~\ref{thm:T3} invokes Lemma~\ref{lem:T1} (TV transport) and Lemma~\ref{lem:T2} (Lyapunov drift). We record an alternative derivation of the key recurrence that applies Proposition~\ref{prop:finite_scale_2adic} directly to the set of survivors, bypassing the TV route entirely. This demonstrates that the geometric decay of the survivor mass does not depend on the TV estimate of Lemma~\ref{lem:T1}, which is vacuous at the working parameters ($\delta_1<0$).

Let $I_k\subset\mathcal B$ be the set of survivors after $k$ blocks, so $|I_k|=Q_k=A_kM$ with $M=2^{\alpha B-1}$. Applying Proposition~\ref{prop:finite_scale_2adic} to $I_k$ with $m=\lceil\sigma d\rceil$:
\[
\frac{|I_{k+1}|}{|I_k|}
\;\le\;2^{-dI_2(\sigma)+O(\log d)}+\frac{\binom{\sigma d}{d}}{|I_k|}
\;=\;c_*+\frac{\binom{\sigma d}{d}}{A_kM}.
\]
Multiplying through by $A_k$ and using $|I_{k+1}|=A_{k+1}M$:
\[
A_{k+1}\;\le\;c_*\,A_k+\frac{\binom{\sigma d}{d}}{M}.\tag{$\star$}
\]
Iterating with $A_0=1$:
\[
A_k\;\le\;c_*^{\,k}+\frac{\binom{\sigma d}{d}}{M\,(1-c_*)}.
\]
The additive term is exponentially small in $B$ by the parameter gap $\delta_{\rm d}:=\alpha-\sigma t\,H_2(1/\sigma)>0$:
\[
\frac{\binom{\sigma d}{d}}{M(1-c_*)}\;\le\;\frac{2^{\,\sigma d\,H_2(1/\sigma)+1}}{2^{\alpha B}(1-c_*)}
\;=\;\frac{2^{\,1-\delta_{\rm d}B}}{1-c_*}.
\]
For $\alpha=1.10$, $\sigma=1.918$, $t=0.30$: $\sigma t\,H_2(1/\sigma)\approx0.5745$, so $\delta_{\rm d}\approx0.5255>0$.

The recurrence $(\star)$ has the same structure as that derived in the proof of Theorem~\ref{thm:T3} via Lemmas~\ref{lem:T1}--\ref{lem:T2}, with the additive term $\binom{\sigma d}{d}/M$ replacing $2^{m+d+1}/M$. Since $\binom{\sigma d}{d}\le2^{\sigma d\,H_2(1/\sigma)}<2^{(\sigma+1)d}\le2^{m+d}$, the alternative additive term is smaller. The Lyapunov drift of Lemma~\ref{lem:T2}(ii) then yields the Hausdorff dimension bound $\dim_H(E_{N_0})\le1+\log_2\rho/(2d)<1$ as before. The role of the present remark is to isolate the TV-free part of the argument: the recurrence $(\star)$, and hence the exponential decay of $A_k$, is a direct consequence of Proposition~\ref{prop:finite_scale_2adic} alone.
\end{remark}

\begin{remark}[Rules honoured; scope]
All TV comparisons are against $\nu_m(r)=2^{-(m-1)}$ on the odd
classes. All rates $c_*=2^{-BtI_2(\sigma)}$ are with $I_2$ in bits. The analytically predicted base survival $c_*^{\rm pred}=2^{-BtI_2(\sigma)} \in[0.52,0.61]$ perfectly matches the measured block survival stationary phase $c_* \in 0.51$--$0.56$ from the numerical observation. Theorem~\ref{thm:T3} is unconditional for each \emph{fixed} $N_0$. It does not imply the full pointwise conjecture: a counterexample lies in $\bigcap_{N_0}E_{N_0}$, a countable intersection of $\dim_H<1$ sets, which may still be non-empty.
\end{remark}

\subsection{Natural-density polylogarithmic descent}

The deterministic $2$-adic transport assembly (first-passage reversal over decreasing thresholds) was recently shown by Shaik (August 2026) to unconditionally bound the measure of exceptional trajectories in natural density. Our exact conditional transport (Lemma~\ref{lem:T1}) natively acts as the strict input for that assembly.

\begin{corollary}[Unconditional polylogarithmic descent in natural density]
For almost all $N$ in natural density, the minimum of the Collatz orbit satisfies
\[
\mathrm{Col}_{\min}(N) \le C (\log N)^A
\]
for any $A > A_{FP} \approx 9.991$, achieved within $c \log N$ unaccelerated steps, with the entire orbit path uniformly bounded by $N^{1+\beta}$ for any $\beta > 0$.
\end{corollary}

\begin{proof}
As demonstrated in recent work [I. A. Shaik, \emph{Polylogarithmic Descent for Almost All Collatz Orbits in Natural Density}, preprint, 2026], assembling the $2$-adic reverse transport along a chain of geometrically decreasing thresholds reduces the union bound of passage times from a linear horizon $O(M)$ to a square-root corridor $O(\sqrt{M \log M})$. Combined with the subcritical binomial timeout tail $\kappa_* = 1 - H_2(\log_3 2) \approx 0.05$ (the exact rate of which we independently verified numerically), this yields the critical exponent $A_{FP} = (2\kappa_*)^{-1}$. 
This route is fully deterministic and independent of the sum-conditioned $3$-adic Fourier decay (Theorem B), closing the almost-all gap unconditionally.
\end{proof}

\begin{remark}[Honest status and synthesis]
The point-wise Collatz conjecture remains open. The contribution of this programme is fivefold: the exact $2$-adic transport theorem, the numerical validation of the $3$-adic pillar, the synthesis exhibiting why the two heads must be treated by different methods, the unconditional Hausdorff dimension drop for the exceptional set at any fixed barrier, and the formal bridging to Shaik's unconditional natural-density polylogarithmic descent. While $2$-adic transport alone determines the bounding exponents, the $3$-adic Fourier spectrum (Allikvere's Theorem B) and our adversarial stress test remain structurally necessary for establishing the \emph{exact distribution} of the first-passage law.
\end{remark}

\section{Research and Software Disclosure}
\label{sec:disclosure}

Generative AI systems (specifically Anthropic Claude/Claude Code, DeepSeek, 
Alibaba Qwen, and Google Gemini) were used throughout this project for the 
following tasks:

\begin{itemize}
\item \textbf{Computational exploration:} Development and debugging of Python 
scripts for parity scanning (\texttt{zone\_parity\_search.py}), confluence 
census (\texttt{confluence\_census.py}), reverse tree generation 
(\texttt{reverse\_tree\_xstar.py}), and Fourier spectral analysis.

\item \textbf{Formal verification:} Lean 4 formalization of structural 
theorems (Proposition B3, Exact Conditional Transport, multi-block survival 
recurrence), including syntax generation, proof skeleton construction, and 
type-checking workflows.

\item \textbf{Diagnostic analysis:} Numerical stress testing of Fourier decay 
hypotheses (Observations~\ref{obs:low_freq_cancellation}, 
\ref{obs:two_regimes}), polymer free energy modeling 
(Observation~\ref{obs:polymer_free_energy}), and adversarial validation of 
sum-conditioned cancellation.

\item \textbf{Structural synthesis:} Organization of the five-result framework 
(F1/R1, S1, I1, C1, M1), identification of closed proof paths 
(Low-Frequency Cancellation, Spectral Transversality, Reverse-chord encoding), 
and formulation of the two-heads dichotomy (Section~\ref{sec:synthesis}).
\end{itemize}

The authors selected all mathematical architectures, verified all theoretical 
claims, performed independent validation of numerical results, and accept full 
responsibility for the content of this paper. No finite computation serves as 
a premise of any theorem; all numerical evidence is explicitly labeled as 
observation or diagnostic.

All code, data, and Lean formalizations are available in the supplementary 
materials and at the project repository.

\section{Conclusion: What Was Mapped, What Was Closed, and What Remains}
\label{sec:conclusion}

This programme set out to understand the Collatz space through the lens of
large-deviation theory, and it ends not with a proof of the conjecture but
with something we believe is equally valuable: a \emph{complete map of the
territory}, with every false path explicitly closed and every open bridge
precisely identified.

\paragraph{The five linked results.}
The theoretical core of the paper is the duality between the \emph{typical}
sector (decaying orbits) and the \emph{rigid} sector (confluence structures),
governed by a single rate function.  Concretely: the forward shift-vector
obeys a Cram\'er LDP with the geometric(2) rate $I_2$ (\textbf{F1}); the
tilted $3$-adic reverse operator satisfies a backward LDP whose normalized
rate coincides with $I_2$ (\textbf{R1}); the generic peak-path is a
\emph{soliton} that avoids all rigid structures (\textbf{S1}); every
confluence center is a \emph{dressed vacuum}---the $(1,1,2)$ fixed point
$\xi=-29/11$ with $O(1)$ total charge (\textbf{I1}); each center sits at the
log-midpoint $\mathrm{bits}(c)=\tfrac12 P+O(1)$ by detailed balance at the
waist (\textbf{C1}); and constructive reverse synthesis of macro-centers is
algorithmically closed by a CRT dimensionality deficit $\Delta\ge 2P$
(\textbf{M1}).  The paradigm crystallises as: \emph{crystals can be found but
not grown.}

\paragraph{The closure of the Fourier route.}
The most consequential negative result of this programme is the systematic
falsification of every candidate Fourier mechanism for the restart/transport
step.  The Low-Frequency Cancellation Lemma (macroscopic peaks $\sim0.6$ at
$M=22,24$ with no decay), Spectral Transversality ($8/10$ unrestricted peaks
aligning with the dominant layer), and the sum-conditioned restart hypothesis
(normalized peaks $=1.0000$, the apparent $O(d^{-1/2})$ decay exposed as a
Gaussian mass artifact) are all conclusively closed
(Observations~\ref{obs:low_freq_cancellation},~\ref{obs:hierarchical_inter_bucket};
Section~\ref{sec:synthesis}).  The two-regime structure of the $3$-adic
Fourier spectrum (fast bulk $L^2$ decay $c_{\mathrm{avg}}\approx0.61$ versus
worst-case $c_\infty\le0.054$ with $c_\infty=0$ not excluded;
Observation~\ref{obs:two_regimes}) and the worst-case barrier
(Theorem~\ref{thm:barrier}, cf.\ Tao Remark~1.18) establish that no $L^2$
averaging can bypass the pointwise Fourier rate.  The polymer free-energy
diagnostic (Observation~\ref{obs:polymer_free_energy}) provides numerical
evidence that $c_\infty>0$ but with a microscopic rate $\sim10^{-6}$,
confirming that the $3$-adic head is theoretically sound yet practically too
slow to bridge the Archimedean gap alone.

\paragraph{The architectural contribution.}
The central structural insight is the \emph{separation of concerns between two
independent heads}.  The $2$-adic head is \emph{exact and cancellation-free}:
the affine word map $T_w$ bijects each cylinder onto $\mathbb{Z}_2$
(Theorem~\ref{thm:exact_conditional_transport}), the conditional endpoint law
on every $S$-layer is exactly Haar, and the $O(1)$ Fourier ``noise'' of the
unrestricted spectrum is entirely the Fourier transform of the initial
Archimedean indicator, not a mixing failure.  The $3$-adic head carries
\emph{all} the cancellation: its sum-conditioned Fourier decay
(Allikvere Theorem~B) is superpolynomial and survives the adversarial stress
test (Observation~\ref{obs:stress_B}), but it is specific to $b=3$ and fails
for $b\ge5$ where the stationary IFS measure is singular.  The threshold
$\delta(b)=\log4/\log b\gtrless1$ is the exact condition for the
Allikvere transport to operate.  This two-heads architecture explains why the
natural-density bridge can stand for $3x+1$ and must collapse for $5x+1$.

\paragraph{New unconditional results.}
Three unconditional results emerge from the synthesis.
\emph{First}, the multi-block survival analysis
(Section~\ref{sec:synthesis}, Theorem~\ref{thm:T3}) establishes
$\dim_H(E_{N_0})\le1+\log_2\rho/(2d)<1$ for each fixed barrier $N_0=2^B$,
with the corrected iteration showing that the Lyapunov residual is a
\emph{constant} resolution floor $C'2^{-\delta B}$, not a growing term; the
Hausdorff bound is taken at the correct depth $k=\Theta(B)$.
\emph{Second}, the exact $2$-adic transport (Lemma~\ref{lem:T1}) provides the
strict input for Shaik's unconditional natural-density assembly, yielding
$\mathrm{Col}_{\min}(N)\le C(\log N)^A$ for $A>A_{\mathrm{FP}}\approx9.991$
in natural density (Corollary in Section~\ref{sec:synthesis}).
\emph{Third}, the family dichotomy $\delta(b)=\log4/\log b$ is validated by
three independent numerical tests (drift, dimension collapse, first-passage
destruction) and identifies $b=3$ as the unique value where the transport
mechanism operates.

\paragraph{What remains open.}
Four bridges remain open, and we state them precisely.
\emph{(i)} The forward restart/transport hypothesis across scales is open and
mechanism-free: every candidate mechanism has been falsified, and the
Archimedean--$2$-adic misalignment---the topological mismatch between
Archimedean sampling intervals and $2$-adic cylinders---is identified as the
true obstruction that logarithmic density bypasses and natural density must
attack via Kronecker orbits and transcendence theory.
\emph{(ii)} The worst-case Fourier rate $c_\infty$ is not determined; if
$c_\infty=0$, Tao's architecture yields a superpolynomial but not an
exponential exceptional-set bound, and the effective threshold $N_0^*$
remains astronomically above the computational frontier $2^{68}$.
\emph{(iii)} The pointwise Collatz conjecture
$\mathrm{Col}_{\min}(N)=1\;\forall N$ is not implied by any density bound:
a counterexample lies in $\bigcap_{N_0}E_{N_0}$, a countable intersection of
$\dim_H<1$ sets that may be non-empty.
\emph{(iv)} The reverse-chord encoding is superseded by an intrinsic entropy
deficit ($h(4/3)\approx1.0817<\log_2 3\approx1.585$), closing the route to
attaching $\gamma_{\mathrm{rev}}\approx0.993$ to the forward exceptional set.

\paragraph{The path forward.}
The map drawn here suggests three directions.
\emph{Non-Fourier transport}: the exact $2$-adic transport theorem
(Theorem~\ref{thm:exact_conditional_transport}) operates without any
Fourier analysis; a restart/transport lemma phrased in the Tree-Wasserstein
metric on the $2$-adic tree, or in entropy rather than Fourier norms,
might bypass the Archimedean wall.
\emph{The $c_\infty$ question}: closing $c_\infty>0$ rigorously (via the
polymer free-energy or a renewal-theoretic argument) would upgrade the
superpolynomial bound to exponential, though the microscopic rate suggests
the effective gain would be modest.
\emph{Pointwise structure}: the dressed-vacuum picture (Theorem~I1) and the
CRT dimensionality obstruction (Theorem~M1) constrain the \emph{shape} of any
hypothetical counterexample---it would have to be a measure-zero rigid object
of $O(1)$ charge, not a statistical accident---but no current method excludes
its existence.

\paragraph{Final word.}
We have not proved the Collatz conjecture.  What we have done is to
distinguish, for the first time in a single coherent framework, the five
different \emph{types} of statement that the word ``Collatz'' conflates:
the typical decay (proven by Tao), the natural-density descent (proven by
Allikvere and Shaik), the Hausdorff dimension of the exceptional set
(established here for fixed barriers), the structural anatomy of the rare
rigid objects (mapped here in detail), and the pointwise conjecture itself
(open since 1937).  The first four are now cleanly separated; the fifth is
seen clearly across the gap.  Whether that gap can be crossed is a question
for mathematics that does not yet exist.

\end{document}
